\documentclass[openany]{book}
\usepackage{lmodern}
\usepackage{amssymb,amsmath}
\usepackage{ifxetex,ifluatex}
\usepackage{fixltx2e} % provides \textsubscript
\ifnum 0\ifxetex 1\fi\ifluatex 1\fi=0 % if pdftex
  \usepackage[T1]{fontenc}
  \usepackage[utf8]{inputenc}
\else % if luatex or xelatex
  \ifxetex
    \usepackage{mathspec}
  \else
    \usepackage{fontspec}
  \fi
  \defaultfontfeatures{Ligatures=TeX,Scale=MatchLowercase}
\fi
% use upquote if available, for straight quotes in verbatim environments
\IfFileExists{upquote.sty}{\usepackage{upquote}}{}
% use microtype if available
\IfFileExists{microtype.sty}{%
\usepackage{microtype}
\UseMicrotypeSet[protrusion]{basicmath} % disable protrusion for tt fonts
}{}
\usepackage[margin=1in]{geometry}
\usepackage{hyperref}
\hypersetup{unicode=true,
            pdftitle={3 LANGKAH KILAT MAHIR ANALISIS DATA},
            pdfauthor={Pribadi Pramudya - pripramudya.com},
            pdfborder={0 0 0},
            breaklinks=true}
\urlstyle{same}  % don't use monospace font for urls
\usepackage{natbib}
\bibliographystyle{apalike}
\usepackage{color}
\usepackage{fancyvrb}
\newcommand{\VerbBar}{|}
\newcommand{\VERB}{\Verb[commandchars=\\\{\}]}
\DefineVerbatimEnvironment{Highlighting}{Verbatim}{commandchars=\\\{\}}
% Add ',fontsize=\small' for more characters per line
\usepackage{framed}
\definecolor{shadecolor}{RGB}{248,248,248}
\newenvironment{Shaded}{\begin{snugshade}}{\end{snugshade}}
\newcommand{\KeywordTok}[1]{\textcolor[rgb]{0.13,0.29,0.53}{\textbf{#1}}}
\newcommand{\DataTypeTok}[1]{\textcolor[rgb]{0.13,0.29,0.53}{#1}}
\newcommand{\DecValTok}[1]{\textcolor[rgb]{0.00,0.00,0.81}{#1}}
\newcommand{\BaseNTok}[1]{\textcolor[rgb]{0.00,0.00,0.81}{#1}}
\newcommand{\FloatTok}[1]{\textcolor[rgb]{0.00,0.00,0.81}{#1}}
\newcommand{\ConstantTok}[1]{\textcolor[rgb]{0.00,0.00,0.00}{#1}}
\newcommand{\CharTok}[1]{\textcolor[rgb]{0.31,0.60,0.02}{#1}}
\newcommand{\SpecialCharTok}[1]{\textcolor[rgb]{0.00,0.00,0.00}{#1}}
\newcommand{\StringTok}[1]{\textcolor[rgb]{0.31,0.60,0.02}{#1}}
\newcommand{\VerbatimStringTok}[1]{\textcolor[rgb]{0.31,0.60,0.02}{#1}}
\newcommand{\SpecialStringTok}[1]{\textcolor[rgb]{0.31,0.60,0.02}{#1}}
\newcommand{\ImportTok}[1]{#1}
\newcommand{\CommentTok}[1]{\textcolor[rgb]{0.56,0.35,0.01}{\textit{#1}}}
\newcommand{\DocumentationTok}[1]{\textcolor[rgb]{0.56,0.35,0.01}{\textbf{\textit{#1}}}}
\newcommand{\AnnotationTok}[1]{\textcolor[rgb]{0.56,0.35,0.01}{\textbf{\textit{#1}}}}
\newcommand{\CommentVarTok}[1]{\textcolor[rgb]{0.56,0.35,0.01}{\textbf{\textit{#1}}}}
\newcommand{\OtherTok}[1]{\textcolor[rgb]{0.56,0.35,0.01}{#1}}
\newcommand{\FunctionTok}[1]{\textcolor[rgb]{0.00,0.00,0.00}{#1}}
\newcommand{\VariableTok}[1]{\textcolor[rgb]{0.00,0.00,0.00}{#1}}
\newcommand{\ControlFlowTok}[1]{\textcolor[rgb]{0.13,0.29,0.53}{\textbf{#1}}}
\newcommand{\OperatorTok}[1]{\textcolor[rgb]{0.81,0.36,0.00}{\textbf{#1}}}
\newcommand{\BuiltInTok}[1]{#1}
\newcommand{\ExtensionTok}[1]{#1}
\newcommand{\PreprocessorTok}[1]{\textcolor[rgb]{0.56,0.35,0.01}{\textit{#1}}}
\newcommand{\AttributeTok}[1]{\textcolor[rgb]{0.77,0.63,0.00}{#1}}
\newcommand{\RegionMarkerTok}[1]{#1}
\newcommand{\InformationTok}[1]{\textcolor[rgb]{0.56,0.35,0.01}{\textbf{\textit{#1}}}}
\newcommand{\WarningTok}[1]{\textcolor[rgb]{0.56,0.35,0.01}{\textbf{\textit{#1}}}}
\newcommand{\AlertTok}[1]{\textcolor[rgb]{0.94,0.16,0.16}{#1}}
\newcommand{\ErrorTok}[1]{\textcolor[rgb]{0.64,0.00,0.00}{\textbf{#1}}}
\newcommand{\NormalTok}[1]{#1}
\usepackage{longtable,booktabs}
\usepackage{graphicx,grffile}
\makeatletter
\def\maxwidth{\ifdim\Gin@nat@width>\linewidth\linewidth\else\Gin@nat@width\fi}
\def\maxheight{\ifdim\Gin@nat@height>\textheight\textheight\else\Gin@nat@height\fi}
\makeatother
% Scale images if necessary, so that they will not overflow the page
% margins by default, and it is still possible to overwrite the defaults
% using explicit options in \includegraphics[width, height, ...]{}
\setkeys{Gin}{width=\maxwidth,height=\maxheight,keepaspectratio}
\IfFileExists{parskip.sty}{%
\usepackage{parskip}
}{% else
\setlength{\parindent}{0pt}
\setlength{\parskip}{6pt plus 2pt minus 1pt}
}
\setlength{\emergencystretch}{3em}  % prevent overfull lines
\providecommand{\tightlist}{%
  \setlength{\itemsep}{0pt}\setlength{\parskip}{0pt}}
\setcounter{secnumdepth}{5}
% Redefines (sub)paragraphs to behave more like sections
\ifx\paragraph\undefined\else
\let\oldparagraph\paragraph
\renewcommand{\paragraph}[1]{\oldparagraph{#1}\mbox{}}
\fi
\ifx\subparagraph\undefined\else
\let\oldsubparagraph\subparagraph
\renewcommand{\subparagraph}[1]{\oldsubparagraph{#1}\mbox{}}
\fi

%%% Use protect on footnotes to avoid problems with footnotes in titles
\let\rmarkdownfootnote\footnote%
\def\footnote{\protect\rmarkdownfootnote}

%%% Change title format to be more compact
\usepackage{titling}

% Create subtitle command for use in maketitle
\newcommand{\subtitle}[1]{
  \posttitle{
    \begin{center}\large#1\end{center}
    }
}

\setlength{\droptitle}{-2em}

  \title{3 LANGKAH KILAT MAHIR ANALISIS DATA}
    \pretitle{\vspace{\droptitle}\centering\huge}
  \posttitle{\par}
    \author{Pribadi Pramudya - pripramudya.com}
    \preauthor{\centering\large\emph}
  \postauthor{\par}
      \predate{\centering\large\emph}
  \postdate{\par}
    \date{2019-01-14}

\usepackage{booktabs}
\renewcommand{\contentsname}{Daftar Isi}
\renewcommand{\chaptername}{Langkah Kilat}

\begin{document}
\maketitle

{
\setcounter{tocdepth}{1}
\tableofcontents
}
\chapter*{Pendahuluan}\label{pendahuluan}
\addcontentsline{toc}{chapter}{Pendahuluan}

\section*{Selamat Datang}\label{selamat-datang}
\addcontentsline{toc}{section}{Selamat Datang}

\begin{center}\includegraphics[width=0.75\linewidth]{images/ebook} \end{center}

Terima kasih sudah bersedia menyempatkan sedikit waktu untuk mengunjungi
website dan membaca buku ini.

Berikut materi yang akan Anda dapatkan:

\begin{itemize}
\tightlist
\item
  \textbf{Kenapa Anda perlu menguasai} keterampilan analisis data
\item
  \textbf{Konsep dasar} analisis data yang tidak diketahui banyak orang
\item
  Satu \textbf{\emph{tool} pengolah data \emph{powerful}} yang perlu
  Anda kuasai
\item
  \textbf{Teknik dasar} analisis data yang perlu Anda pelajari
\item
  \textbf{Panduan praktis dan rinci} cara menganalisis data dari nol
  sama sekali
\end{itemize}

Versi pdf buku elektronik pendek 3 Langkah Kilat Mahir Analisis Data
bisa diunduh dengan mendaftarkan email Anda terlebih dulu di
\url{https://pripramudya.com}.

Buku dan website ini tersedia secara bebas, di bawah
\href{https://creativecommons.org/licenses/by/4.0/deed.id}{Lisensi CC BY
4.0}.

Anda diperbolehkan:

\textbf{Berbagi} --- menyalin dan menyebarluaskan kembali materi ini
dalam bentuk atau format apapun;

\textbf{Adaptasi} --- menggubah, mengubah, dan membuat turunan dari
materi ini untuk kepentingan apapun, termasuk kepentingan komersial.

Dengan \textbf{mencantumkan sumber tautan} asli dari konten terkait.

\section*{Kenapa Anda Perlu Menguasai Keterampilan Analisis
Data}\label{kenapa-anda-perlu-menguasai-keterampilan-analisis-data}
\addcontentsline{toc}{section}{Kenapa Anda Perlu Menguasai Keterampilan
Analisis Data}

Data ada di mana-mana.

Hampir setiap hari Anda berhadapan dengan data, apapun profesi Anda.

Apalagi di zaman \emph{now}, data berkembang begitu pesat dalam hal
jumlah \emph{(volume)}, variasi \emph{(variety)}, dan kecepatan
\emph{(velocity)}.

Ya, sekarang kita berada di era \textbf{big data}.

Begitu melimpah ruahnya, sehingga tidak banyak orang yang memiliki
keterampilan mumpuni untuk mengambil \emph{insight} (wawasan baru) dari
banyaknya data tersebut.

Tidak perlu jauh mencari sumber data dari luar, data yang tersedia di
dalam bidang keprofesian Anda pun sepertinya masih belum dapat digali
dengan optimal.

Di tempat kerja saya, pegawai yang menguasai keterampilan analisis data
memiliki peluang lebih baik dalam menanjak karir. Apalagi jika hasil
analisis data ini digunakan sebagai bahan presentasi oleh kepala divisi.
Karirnya meroket.

Ini salah satu contoh nyata. Manajemen begitu membutuhkan \emph{insight}
yang dapat memberikan suatu perubahan lebih baik dalam mengelola roda
aktivitas perusahaan sehingga berujung pada peningkatan profitabilitas
perusahaan. Manajemen tingkat atas pasti tidak akan ragu memberikan
promosi bagi pegawai yang memberikan \emph{insight} tersebut. Peluang
\emph{(opportunity)} penanjakan karir pegawai yang memiliki keterampilan
analisis yang mumpuni lebih baik dari pegawai yang memiliki keterampilan
ala kadarnya saja.

Saya akan meringkas bab ini cukup dengan dua kata saja:

\emph{\textbf{INSIGHT \& OPPORTUNITY}}

Ingin mendapatkan \emph{insight} yang dapat menjawab permasalahan di
bidang keprofesian Anda dan membuka \emph{opportunity} karir terbaik
Anda?

\textbf{Kuasai keterampilan analisis data.}

\section*{Konsep Dasar Analisis Data}\label{konsep-dasar-analisis-data}
\addcontentsline{toc}{section}{Konsep Dasar Analisis Data}

``Saya sudah tahu apa itu analisis data, tidak perlu dijelaskan lagi.''

Mungkin hal tersebut yang ada di benak Anda ketika membaca judul di bab
ini.

Memang, membaca teori tidaklah begitu menarik. Anda ingin langsung
meloncat ke bab langkah kilat.

Silakan saja. Gunakan menu di daftar isi lalu klik ke bab yang ingin
Anda baca.

Namun, tetap saya sarankan untuk membaca bab ini sampai akhir.

Anda perlu tahu konsep dasar analisis data agar memudahkan Anda agar
bisa cepat menguasai keterampilan ini.

\section*{Apa itu Analisis Data}\label{apa-itu-analisis-data}
\addcontentsline{toc}{section}{Apa itu Analisis Data}

Jadi, apa sebenarnya analisis data itu?

Analisis data adalah proses mengetahui, memahami, memilah, memecah,
merinci, mengurai data dengan berbagai cara atau teknik tertentu seperti
manipulasi, visualisasi, pengelompokkan, perbandingan, atau teknik
lainnya sehingga didapatkan informasi berguna yang dapat menjawab
permasalahan serta membantu dalam pengambilan keputusan.

Ya, \textbf{tujuan dari analisis data adalah pengambilan keputusan yang
lebih baik}. Keputusan yang didasarkan pada informasi hasil analisis
data bukan berdasarkan naluri semata. Keputusan luar biasa krusialnya
yang bisa saja membuat sebuah perusahaan sebelumnya berkibar berujung
menjadi gulung tikar ataupun sebaliknya.

Ini membuat analisis data menjadi salah satu keterampilan yang sangat
penting penerapannya di hampir seluruh bidang.

Beberapa contoh penerapan dalam bidang bisnis sebagai berikut:

\textbf{Perkenalan produk kopi terbaru Starbucks}

Starbucks memperkenalkan produk kopi terbaru dengan memperhatikan apakah
pelanggan menyukai produk baru tersebut. Pada pagi hari produk tersebut
mulai dijual, Starbucks memantau berbagai blog, Twitter, forum diskusi
dan grup tentang kopi untuk mengetahui bagaimana reaksi pelanggan. Siang
harinya, Starbucks menemukan bahwa walaupun pelanggan menyukai rasa
kopinya, namun mereka berpikir harganya terlalu mahal. Starbucks
langsung menurunkan harga. Pada sore harinya, semua persepsi negatif
terkait produk kopi baru tersebut hilang.

Pendekatan dengan respons cepat seperti ini tentu lebih baik
dibandingkan pendekatan tradisional dengan menunggu laporan penjualan
harian masuk lebih dulu yang pada akhirnya ditemukan bahwa penjualannya
mengecewakan. Langkah \emph{jadul} berikutnya mencari tahu alasan kenapa
penjualan mengecewakan dengan melakukan diskusi grup terpusat
\emph{(focus group discussion)}. Selanjutnya, beberapa minggu kemudian
Starbucks baru menemukan bahwa harganya terlalu mahal lalu menurunkan
harganya.

\textbf{Sistem rekomendasi Amazon}

Amazon menggunakan data penjualan, data produk yang sering dilihat, dan
perilaku serta preferensi pelanggan lainnya untuk membuat sebuah sistem
rekomendasi. Sistem ini membantu Amazon meningkatkan 29\% penjualannya.
Jika Anda mengunjungi situs Amazon, di bagian bawah Anda akan melihat
beberapa produk rekomendasi, produk yang Anda lihat sebelumnya, sampai
dengan produk terlaris serta tawaran produk menarik lainnya.

\textbf{Efisiensi Chevron dalam aktivitas pengeboran}

Biaya setiap pengeboran di selat Meksiko diperkirakan sampai dengan
\$100 juta. Rugi sekali jika hasil pengeboran tersebut tidak didapat
sumber minyak bumi. Untuk meningkatkan peluang keberhasilan menemukan
sumber minyak, Chevron menganalisis 50 terabyte data seismik. Dengan
dukungan kecanggihan komputer, kapasitas penyimpanan, serta tim analis
yang mumpuni berhasil meningkatkan keberhasilan menemukan sumber minyak
dari yang sebelumnya 1 dari 5 pengeboran menjadi 1 berbanding 3.

\section*{Proses dan Tahapan Analisis
Data}\label{proses-dan-tahapan-analisis-data}
\addcontentsline{toc}{section}{Proses dan Tahapan Analisis Data}

Untuk menghasilkan informasi dan \emph{insight} yang dapat memberikan
keputusan lebih baik, Anda harus tahu proses dan tahapan dari analisis
data. Tahapan analisis data terdiri dari 6 tahapan.

\begin{enumerate}
\def\labelenumi{\arabic{enumi}.}
\tightlist
\item
  Menentukan Tujuan
\item
  Menentukan Metode Pengukuran yang Digunakan
\item
  Pengumpulan Data
\item
  Pembersihan dan Pembentukan Data
\item
  Analisis Data
\item
  Interpretasi dan Komunikasi Hasil Analisis Data
\end{enumerate}

\subsection*{1. Menentukan Tujuan}\label{menentukan-tujuan}
\addcontentsline{toc}{subsection}{1. Menentukan Tujuan}

Ini adalah tahap pertama dari proses analisis data. Di bidang
keprofesian Anda, Anda harus menentukan tujuan dari permasalahan yang
akan diselesaikan dimulai dengan pertanyaan yang tepat. Tahap ini cukup
vital mengingat pertanyaan yang salah akan menghasilkan informasi atau
\emph{insight} yang tidak tepat pula.

``Kenapa total penjualan bulan ini turun dibandingkan bulan lalu?''

``Kenapa 7 dari 10 proyek yang sedang berjalan saat ini terlambat?''

``Bagaimana cara mencegah pelanggan kita pindah ke kompetitor?''

Itu adalah contoh permasalahan yang akan disolusikan.

Dengan menentukan pertanyaan yang tepat, Anda bisa lebih fokus dalam
mencapai tujuan serta lebih efektif untuk menyelesaikan permasalahan
yang ada.

Sumber data yang akan digunakan nantinya juga tergantung dari penentuan
tujuan dan pertanyaan tepat dari tahap ini.

\subsection*{2. Menentukan Metode Pengukuran yang
Digunakan}\label{menentukan-metode-pengukuran-yang-digunakan}
\addcontentsline{toc}{subsection}{2. Menentukan Metode Pengukuran yang
Digunakan}

Anda sudah meneentukan tujuan di tahap pertama, selanjutnya adalah apa
yang perlu dijadikan tolak ukur dan bagaimana cara Anda mengukurnya.

Misalnya Anda mengambil contoh:

``Kenapa total penjualan bulan ini turun dibandingkan bulan lalu?''

Dari permasalahan ini, Anda bisa tentukan bahwa yang harus diukur adalah
\textbf{penjualan}. Lalu, Anda tentukan lagi metrik dari penjualan,
\textbf{jumlah barang} atau \textbf{jumlah uang} atau keduanya.

\begin{itemize}
\tightlist
\item
  Metrik jumlah barang: buah, lusin, kodi, dus, porsi, dan lainnya.
\item
  Metrik jumlah uang: ratusan ribu rupiah, jutaan rupiah, ribuan dollar,
  dan lainnya.
\end{itemize}

Pengukuran juga bisa dibandingkan berdasarkan periodenya. Terhadap bulan
sebelumnya atau terhadap tahun sebelumnya di bulan yang sama.

Selain itu, Anda pikirkan juga kira-kira faktor penyebab permasalahan
yang bisa dijadikan sebagai tambahan pengukuran. Contoh di bidang
penjualan, faktor seperti jumlah pembeli, promosi, diskon, lokasi,
harga, bisa dimasukkan sebagai tambahan metode pengukuran.

Pastikan pengukuran ini dilakukan secara tepat karena nantinya akan
berpengaruh terhadap kualitas hasil analisis dan pengambilan keputusan.

\subsection*{3. Pengumpulan Data}\label{pengumpulan-data}
\addcontentsline{toc}{subsection}{3. Pengumpulan Data}

Anda sudah menentukan tujuan dan juga metode pengukuran di tahapan
sebelumnya. Selanjutnya adalah mencari dan mengumpulkan data yang
relevan. Sumber data bisa dari mana saja. Internet, basis data internal,
survey, wawancara, dan lainnya.

Contoh permasalahan penjualan di tahapan sebelumnya. Data penjualan,
promosi, jumlah pelanggan, bisa didapat dari internal. Untuk survey
kepuasan pelanggan bisa dilakukan dengan metode pengisian kuesioner
ataupun dengan wawancara langsung.

Yang perlu diingat adalah Anda harus mengumpulkan berbagai macam data
tersebut secara terstruktur sehingga bisa langsung dapat diolah.

\subsection*{4. Pembersihan dan Pembentukan
Data}\label{pembersihan-dan-pembentukan-data}
\addcontentsline{toc}{subsection}{4. Pembersihan dan Pembentukan Data}

Dari seluruh tahapan proses analisis data, tahap ini adalah tahap yang
paling banyak menyita waktu Anda. Ketika Anda membuka data yang Anda
kumpulkan di tahap 3, Anda menemukan struktur data yang tidak siap olah.
Berantakan. Duplikasi, salah ketik, karakter spesial, huruf
kapital-kecil, data kosong atau NA, format tanggal/waktu tidak sama, dan
lainnya merupakan beberapa hal yang sering ditemui pada sebuah basis
data.

Untuk mempermudah pembersihan dan pembentukan data ini Anda bisa
menggunakan \emph{tools} sejuta umat, Microsoft Excel. Anda juga bisa
menggunakan software pengolah data lain yang Anda kuasai, misalnya SPSS,
SAS, STATA, Python, R atau lainnya.

Jangan terlalu menganggap remeh tahapan ini. Buat dan bentuk data sebaik
mungkin sehingga data bisa langsung dianalisis. Jika data tidak bersih,
semahir apapun Anda dalam menganalisis data, hasilnya tidak akan
optimal. Ingat, jika masukannya sampah, keluarannya juga sampah.

\subsection*{5. Analisis Data}\label{analisis-data}
\addcontentsline{toc}{subsection}{5. Analisis Data}

Setelah data dibersihkan dan dibentuk, tahapan penting berikutnya adalah
tahap analisis data. Tujuan dari tahapan ini adalah untuk mengerti lebih
dalam tentang data beserta variabel-variabelnya. Ada beberapa teknik
dasar analisis data yang dapat Anda gunakan.

\begin{itemize}
\tightlist
\item
  Teknik Eksplorasi
\item
  Teknik Visualisasi
\end{itemize}

\textbf{Teknik Eksplorasi}

Sesuai dengan namanya, teknik menjelajah untuk mengerti tentang data.
Anda memulainya dengan mencari tahu berapa jumlah baris dan kolom dari
basis data, lalu melihat jenis variabelnya: karakter, numerik, atau
kategorikal. Dilanjut dengan meringkas data tersebut sehingga dapat
menampilkan informasi penting seperti variabel atau kategori dengan
frekuensi terbanyak, nilai tertinggi, rerata, dan informasi lainnya.

Intinya Anda menggunakan fungsi-fungsi dari \emph{tools} analisis untuk
mengerti tentang data. Terkadang Anda perlu kembali melakukan proses di
tahapan sebelumnya, yakni pembersihan data lalu balik lagi ke tahap ini.
Ya, tahap 4 dan 5 merupakan proses yang iteratif. Hal ini dilakukan agar
Anda bisa lebih cepat mengerti tentang informasi yang tersimpan dalam
data.

\textbf{Teknik Visualisasi}

\begin{quote}
``Sebuah gambar bernilai ribuan kata.''
\end{quote}

Salah satu pepatah bahasa inggris ini benar adanya. Anda lebih cepat
mengerti data dengan mereprentasikannya dalam bentuk grafik seperti
diagram batang, diagram garis, diagram tebar, histogram, serta info
grafik lainnya. Teknik visualisasi merupakan teknik termudah, tercepat,
dan terefektif dalam menampilkan informasi tentang data.

Selain itu, hasil dari teknik visualisasi ini bisa digunakan sebagai
alat bantu dalam mengkomunikasikan informasi dan insight.

Jika ada satu teknik analisis data yang saya rekomendasikan untuk
dikuasai pertama kali, maka teknik visualisasi inilah yang perlu Anda
pelajari secepatnya.

\subsection*{6. Interpretasi dan Komunikasi Hasil Analisis
Data}\label{interpretasi-dan-komunikasi-hasil-analisis-data}
\addcontentsline{toc}{subsection}{6. Interpretasi dan Komunikasi Hasil
Analisis Data}

Hasil analisis yang sudah Anda selesaikan harus diinterpretasikan.
Maksud dari interpretasi ini adalah mengubah hasil analisis yang masih
dalam bentuk teknikal menjadi temuan, informasi, ataupun \emph{insight}
yang bisa dimengerti oleh orang awam sekalipun.

Selanjutnya, apakah hasil interpretasi ini sudah menjawab pertanyaan
permasalahan di proses tahap 1.

Buat kesimpulan, akan lebih baik lagi jika disertai rekomendasi dan
langkah-langkah yang diperlukan untuk menyelesaikan permasalahan yang
ada.

Sebaik apapun Anda dalam menganalisis data itu semua akan percuma jika
Anda tidak mengkomunikasikan rekomendasi, informasi, dan insight ke
atasan, manajemen, klien, pelanggan, masyarakat, ataupun
\emph{stakeholder} Anda. Bukan berarti teknik analisis data tidak
penting sama sekali, namun dengan interpretasi dan mengkomunikasikan
hasil analisis maka dampak efektivitas dalam mensolusikan permasalahan
akan lebih terasa.

Media dalam mengkomunikasikan pun terdapat banyak alternatif. Anda bisa
membuat hasil analisis data dalam bentuk infografik (file gambar),
tulisan dalam format pdf, slide presentasi, \emph{spreadsheet},
dahsboard, sampai dengan artikel di blog.

Dengan mengkomunikasikan hasil analisis data, maka pintu
\emph{opportunites} terbuka lebih lebar, mengundang karir terbaik Anda.

\section*{3 Langkah Kilat Mahir Analisis
Data}\label{langkah-kilat-mahir-analisis-data}
\addcontentsline{toc}{section}{3 Langkah Kilat Mahir Analisis Data}

Anda sudah mengetahui tentang konsep dasar data analisis yang menjadi
landasan dan pondasi Anda dalam mempelajari keterampilan analisis data
pada bab sebelumnya. Konsep tersebut merupakan hal yang penting agar
Anda dapat cepat mahir dan menguasai keterampilan analisis data.

Di bab selanjutnya akan dibahas bagaimana mempelajari dan
mengaplikasikan teknik analisis data. Anda akan dipandu secara rinci
bagaimana agar Anda dapat dengan cepat menguasai keterampilan analisis
data dengan 3 langkah kilat berikut.

\chapter{\texorpdfstring{Instal dan Pelajari Satu \emph{Tool} Analisis
Data yang \emph{Powerful}
ini}{Instal dan Pelajari Satu Tool Analisis Data yang Powerful ini}}\label{instal-dan-pelajari-satu-tool-analisis-data-yang-powerful-ini}

Salah satu yang mempengaruhi proses dan kualitas hasil analisis adalah
dari \emph{tool} yang dipakai. \emph{Tool} atau perangkat lunak ini
memudahkan Anda sebagai analis untuk melakukan proses analisis data
secara efektif dan efisien. Selain itu, jika Anda menguasai dan mahir
salah satu \emph{tool} analisis data maka peluang dalam membuka karir
terbaik Anda akan lebih tinggi. Begitu pentingnya pemilihan \emph{tool}
ini, sehingga saya akan merekomendasikan terlebih dahulu beberapa
alternatif \emph{tool} analisis data yang dipakai banyak orang.

\section{\texorpdfstring{\emph{Tool} atau Perangkat Lunak Analisis Data
Populer yang Banyak Orang
Pakai}{Tool atau Perangkat Lunak Analisis Data Populer yang Banyak Orang Pakai}}\label{tool-atau-perangkat-lunak-analisis-data-populer-yang-banyak-orang-pakai}

Ada 3 \emph{tool} yang akan dibahas di bagian ini:

\begin{itemize}
\tightlist
\item
  Microsoft Excel
\item
  Python
\item
  R
\end{itemize}

\subsection{Microsoft Excel}\label{microsoft-excel}

\emph{Software} nomor satu yang paling banyak dipakai orang di dunia.

Microsoft pertama kali membuat dan memasarkan sebuah \emph{spreadsheet
program} bernama Multiplan di tahun 1982 yang sangat populer di CP/M
sistem (sistem operasi mikrokontroler), namun kalah bersaing dengan
Lotus 1-2-3 di MS-DOS. Hal ini membuat Microsoft mengembangkan versi
terbaru dari \emph{spreadsheet program} bernama Excel yang dimaksudkan
untuk menyaingi seluruh fitur Lotus 1-2-3 dan membuatnya lebih baik.

Versi pertama Excel dipasarkan untuk Mac di tahun 1985 dan pada Nopember
1987 untuk Windows. Di tahun 1988, Excel berhasil memimpin pangsa pasar
\emph{spreadsheet program} menyingkirkan Lotus 1-2-3 dan menjadi
\emph{software} yang paling banyak dipakai di seluruh dunia.

Yang pasti Anda sudah familiar dan mahir menggunakan \emph{sofware} ini.
Jika Anda sedang mencari pekerjaan kantoran, maka keterampilan dasar
yang dijadikan sebagai syarat di sebagian besar perusahaan adalah bisa
menggunakan \emph{sofware} Excel. Saking mudahnya mengoperasikan Excel,
pengguna pemula tidak perlu belajar khusus untuk menguasainya. Tidak
perlu tahu tentang \emph{programming}. Cukup bermodalkan \emph{point,
click,} dan \emph{drag}, Anda sudah bisa melakukan analisis data dengan
\emph{software} ini.

\subsection{Python}\label{python}

Python adalah salah satu bahasa pemrograman untuk membuat berbagai macam
aplikasi \emph{(general purpose programming)}. Dibuat oleh Guido van
Rossum dan pertama kali dirilis pada tahun 1991, Python menjadi salah
satu bahasa pemrograman yang banyak digunakan di dunia. Kemudahan
\emph{syntax}, lisensi \emph{open source}, serta dukungan komunitas
merupakan beberapa faktor yang membuat popularitas Python melesat di
beberapa tahun belakangan ini.

Python juga menjadi salah satu \emph{tool} yang sering digunakan untuk
analisis data. Dukungan dari banyaknya paket membuat Python sangat
fleksibel. Ada NumPy dan Pandas untuk manipulasi data, SciPy untuk
perhitungan \emph{science}, Matplotlib untuk data visualisasi, serta
paket-paket lainnya.

IDE \emph{(Integrated Development Environment)} Python juga ada banyak
pilihan. Mulai dari Spyder, Jupyter Notebook, Atom, PyCharm, dan Rodeo.
Banyaknya pilihan dan fleksibilitas serta kemudahan dalam programming
membuat Python menjadi \emph{tool} nomor satu untuk data analisis saat
ini.

\subsection{R}\label{r}

R adalah perangkat lunak gratis untuk komputasi dan analisis data
statistik yang juga dapat menampilkan visualisasi grafik. R didukung
oleh ribuan \emph{packages} yang dibuat oleh banyak kontributor aktif
sehingga memudahkan pengguna dalam pengembangan aplikasi.

R pertama kali dibuat oleh Ross Ihaka and Robert Gentleman di University
of Auckland, New Zealand pada tahun 1992. Saat ini dikembangkan oleh R
Development Core Team. Dinamakan R diambil dari abjad pertama kedua
pembuat program ini. Versi pertama dirilis pada tahun 1995 dan versi
beta yang stabil pada tahun 2000.

Seperti halnya Python, bahasa pemrograman memiliki lisensi \emph{open
source}, dan dapat diinstal di seluruh sistem operasi: Windows, Linux,
ataupun Mac. Komunitasnya pun banyak ditemukan di website seperti
Stackoverflow ataupun di Github.

Terkait popularitas, R masih kalah dibandingkan dengan Python. R nomor
dua dalam hal \emph{tool} analisis data paling banyak digunakan dan
paling populer setelah Python.

Anda masih ingat judul bab ini? Tentang \emph{tool} analisis data
\emph{powerful}.

Ya, \textbf{R adalah \emph{tool} analisis data \emph{powerful} yang
perlu Anda pelajari dan kuasai.}

\section{Kenapa harus R}\label{kenapa-harus-r}

Anda saat ini pasti bertanya-tanya kenapa harus belajar R? Kenapa tidak
Excel saja atau Python?

Yang pertama, kenapa tidak Excel saja? Karena Anda pasti sudah cukup
mahir menggunakannya. Apakah hanya Anda yang mahir Excel? Tentu saja,
tidak hanya Anda yang mahir, tetapi juga banyak orang. Menurut saya,
orang yang minimal berpendidikan SMA pasti mengenal Excel.

Jika Anda dalam posisi bersaing untuk mendapatkan karir yang lebih baik
di suatu perusahaan atau mencari pekerjaan baru di perusahaan lain, maka
Anda perlu memiliki \emph{skill} yang berbeda dibandingkan kebanyakan
orang.

\textbf{R atau Python dapat memberikan perbedaan itu.}

Yang kedua, kenapa tidak Python?

\emph{Karena saya belum pernah pakai Python, he he :D}

Pada saat saya bertekad untuk belajar analisis data, saya memilih R
sebagai \emph{tool} utama. Python dan R sama-sama powerful, tinggal
preferensi masing-masing saja. Jika Anda memiliki latar belakang
\emph{software developer}, mungkin Anda lebih cocok dengan Python,
begitu pun sebaliknya. Saya memilih R karena saya bukan \emph{software
developer}. R sangat membantu meningkatkan produktifitas dalam
menganalisis data di tempat kerja saya. R mungkin cocok buat Anda yang
tidak memiliki latar belakang di bidang komputer atau \emph{software}.

Setidaknya ada 3 alasan kenapa Anda harus mencoba belajar R:

\subsection{Dukungan Paket R}\label{dukungan-paket-r}

Jumlah paket R yang bisa didownload sampai dengan sekarang ini sebanyak
13.409 di \href{https://cran.r-project.org/}{CRAN} (Comprehensive R
Archive Network), \emph{repository} R. Paket tersedia sebanyak ini belum
termasuk ribuan lainnya yang ada di GitHub, juga paket komersial yang
dikembangkan oleh vendor besar seperti Microsoft dan Oracle.

Kenapa bisa begitu banyak paket?

\begin{itemize}
\item
  R memiliki banyak komunitas pengembang sangat aktif yang berkontribusi
  dalam membuat paket-paket baru ke CRAN setiap hari. Hal ini membuat
  hampir segala analisis yang ingin Anda lakukan sudah tersedia dalam
  bentuk paket di CRAN.
\item
  R dibuat sejak tahun 1992, yang berarti sampai dengan saat ini sudah
  berusia lebih dari 20 tahun. Selama itu pula, berbagai macam paket
  terus dikembangkan sehingga R menjadi salah satu \emph{tool} analisis
  data yang memiliki banyak pilihan paket dibandingkan \emph{tool}
  analisis data lainnya.
\item
  R banyak digunakan sebagai \emph{tool} utama untuk penelitan
  statistik. Ketika metode baru dikembangkan, hasil penelitian tersebut
  tidak hanya dipublikasikan dalam bentuk makalah atau jurnal tetapi
  juga dalam bentuk paket. Ini membuat R selalu terdepan dalam
  mengadopsi metodologi terbaru.
\item
  R didesain sebagai bahasa antar muka bahasa pemrograman lainnya.
  Banyak paket berfungsi sebagai pendukung perangkat lunak \emph{open
  source} lainnya sehingga membuat R menjadi \emph{tool} yang
  mempermudah kolaborasi berbagai macam metode dan algoritma
  pemrograman.
\item
  Sistem CRAN merupakan platform yang sangat efektif dalam proses
  kolaborasi pengembangan paket R. Dengan sistem yang sudah teratur,
  pengguna bisa dengan mudah membuat, mengembangkan, mengetes, dan
  mendistribusikan paket.
\end{itemize}

\subsection{Belajar Analisis Data lebih mudah menggunakan
R}\label{belajar-analisis-data-lebih-mudah-menggunakan-r}

R menyediakan paket pembelajaran yang dinamakan
\href{https://swirlstats.com/}{swirl}. Swirl mengajarkan Anda
pemgrograman R dan analisis data secara interaktif langsung di dalam
aplikasi R sesuai kecepatan pemahaman Anda.

Selain adanya paket belajar mandiri di swirl, R juga memiliki syntax
atau perintah yang lebih simpel dan lebih mudah diingat sehingga dapat
mempercepat proses pembelajaran Anda. Hal ini juga membuat Anda akan
lebih mudah mengerti script pemrograman yang ditulis oleh orang lain
tanpa perlu membaca komentar programnya.

Ditambah dengan paket-paket yang mempermudah Anda dalam melakukan
analisis data seperti dplyr (untuk manipulasi data), ggplot
(visualisasi), dan lainnya. Paket tersebut dibuat sedemikian rupa
sehingga Anda dibantu untuk mengerti proses dan teknik, bukan syntax
sehingga Anda akan lebih diarahkan pada pembelajaran konsep.

\subsection{Fleksibilitas R}\label{fleksibilitas-r}

\href{https://www.rstudio.com/}{RStudio} adalah salah satu IDE
(Integrated Development Environment) pemrograman R yang sangat populer.
Dengan RStudio, membuat dan mengembangkan aplikasi R menjadi lebih mudah
dibandingkan menggunakan konsol original R.

RStudio sangat powerful dan membuat R menjadi sangat fleksibel. Anda
bisa melakukan banyak hal melalui RStudio. Mulai dari membuat
pemrograman script, grafik, laporan dalam bentuk pdf, ppt, excel, bahkan
ebook, sampai dengan administrasi blog.

Ya, ebook ini dan situs \url{pripramudya.com} dibuat menggunakan
RStudio.

Begitulah cerita tentang R. Jika Anda yakin untuk mulai belajar R,
jangan ragu lagi untuk melakukan langkah kecil berikutnya.

\section{Instal R dan RStudio}\label{instal-r-dan-rstudio}

Pilihlah paket instalasi sesuai sistem operasi yang Anda pakai. Untuk
RStudio, pilih yang versi desktop.

\begin{itemize}
\tightlist
\item
  R dapat diunduh di \href{https://repo.bppt.go.id/cran/}{sini}.
\item
  RStudio dapat diunduh di
  \href{https://www.rstudio.com/products/rstudio/download/}{sini}.
\end{itemize}

Buka software RStudio.

Layout RStudio terdiri dari bagian Script, Console, Variable, dan
Output.

\includegraphics[width=1\linewidth]{images/Rstudio_Layout}

\section{Instal Paket R}\label{instal-paket-r}

Paket adalah komponen penting yang memungkinkan pengguna untuk
menggunakan program siap pakai tanpa harus membuatnya dari awal. Ada
banyak ribuan paket di CRAN, tapi untuk saat ini, saya sarankan untuk
instal cukup satu paket analisis data yang paling powerful, yakni
\textbf{tidyverse}.

Anda dapat menginstalnya cukup dengan mengetik di CONSOLE menggunakan
perintah:

\begin{Shaded}
\begin{Highlighting}[]
\KeywordTok{install.packages}\NormalTok{(}\StringTok{"tidyverse"}\NormalTok{)}
\end{Highlighting}
\end{Shaded}

Cukup instal paket tidyverse ini, Anda sudah bisa mulai belajar analisis
data menggunakan R.

\section{\texorpdfstring{\emph{Syntax} dan Operasi Dasar
R}{Syntax dan Operasi Dasar R}}\label{syntax-dan-operasi-dasar-r}

Di bagian ini, Anda akan diarahkan untuk mengenal \emph{syntax} serta
operasi dasar dari pemrograman R. Apa itu \emph{syntax}?

\emph{Syntax} adalah aturan menulis kata perintah di CONSOLE agar dapat
dimengerti dengan benar oleh bahasa pemrograman R. Aturan \emph{syntax}
ini secara mutlak harus dipenuhi. Jika ada kesalahan penulisan
\emph{syntax}, maka akan ada umpan balik berupa pesan eror yang diterima
oleh pengguna.

Jadi, bagian ini penting untuk Anda ketahui. Tutorial \emph{syntax} dan
operasi dasar ini harus Anda ketik di CONSOLE RStudio.

\hypertarget{kalkulator-interaktif}{\subsection{Kalkulator
Interaktif}\label{kalkulator-interaktif}}

R bisa digunakan sebagai kalkulator. Ketik \texttt{2\ +\ 5} di CONSOLE
dan tekan Enter. Anda akan mendapat umpan balik berupa hasil perhitungan
yang telah Anda input sebelumnya.

\begin{Shaded}
\begin{Highlighting}[]
\DecValTok{2} \OperatorTok{+}\StringTok{ }\DecValTok{5}
\end{Highlighting}
\end{Shaded}

\begin{verbatim}
## [1] 7
\end{verbatim}

Anda juga bisa menggunakan operator aritmatika lainnya seperti
\texttt{*} untuk perkalian, \texttt{-} untuk pengurangan, \texttt{/}
untuk bagi, ataupun \texttt{\^{}} untuk kuadrat.

\subsection{Variabel}\label{variabel}

Jika Anda ingin melakukan perhitungan kedua menggunakan hasil di atas,
maka Anda dapat menyimpannya dalam bentuk variabel. Dengan variabel,
Anda tidak perlu ketik ulang \texttt{2\ +\ 5} setiap kali Anda
memerlukannya.

Untuk membuat variabel baru, gunakan operator \texttt{panah\ ke\ kiri},
seperti ini: \texttt{\textless{}-}. Operator variabel ini merupakan
gabungan karakter `kurang dari' yang diikuti karakter `minus'.

Kenapa \texttt{panah\ ke\ kiri}? Operator ini merupakan simbol yang Anda
gunakan seolah Anda memasukkan nilai yang ada di sebelah kanan panah ke
variabel yang ada di sebelah kiri panah.

Contohnya, Anda akan menyimpan hasil dari \texttt{2\ +\ 5} ke variabel
bernama \texttt{x}. Ketik \texttt{x\ \textless{}-\ 2\ +\ 5} di CONSOLE
dan tekan Enter.

\begin{Shaded}
\begin{Highlighting}[]
\NormalTok{x <-}\StringTok{ }\DecValTok{2} \OperatorTok{+}\StringTok{ }\DecValTok{5}
\end{Highlighting}
\end{Shaded}

Anda akan menyadari, R tidak langsung mengirim umpan balik hasil dari
kalkulasi tersebut. Yang R lakukan hanya menyimpan hasil perhitungan
tersebut ke variabel `x' dan mengasumsikan Anda akan menggunakan varibel
ini nantinya. Daftar nama variabel yang Anda buat akan tampil pada
jendela VARIABLE di RStudio.

Jika Anda ingin menampilkan isi variabel `x', Anda ketik \texttt{x} dan
tekan Enter.

\begin{Shaded}
\begin{Highlighting}[]
\NormalTok{x}
\end{Highlighting}
\end{Shaded}

\begin{verbatim}
## [1] 7
\end{verbatim}

Sekarang, simpan hasil dari \texttt{x\ *\ 3} pada variabel bernama `y'.

\begin{Shaded}
\begin{Highlighting}[]
\NormalTok{y <-}\StringTok{ }\NormalTok{x }\OperatorTok{*}\StringTok{ }\DecValTok{3}
\end{Highlighting}
\end{Shaded}

Tampilkan hasil dari `y'.

\begin{Shaded}
\begin{Highlighting}[]
\NormalTok{y}
\end{Highlighting}
\end{Shaded}

\begin{verbatim}
## [1] 21
\end{verbatim}

\subsection{Vektor}\label{vektor}

Kumpulan dari beberapa objek disebut vektor. Objek bisa berupa angka,
karakter, string, dan lainnya. Vektor berupa angka merupakan bentuk
simpel dari struktur data di R. Bahkan, 1 angka bisa dikatakan sebagai
vektor dengan elemen sepanjang 1.

Cara membuat sebuah vektor adalah dengan menggunakan fungsi \texttt{c()}
yang merupakan kepanjangan dari `concatenate' atau `combine'.

Anda akan membuat sebuah vektor yang berisi nilai numerik 1.1, 9, dan
3.14 (pada penulisan angka, tanda titik di pengaturan standar RStudio
merupakan tanda koma, begitu juga sebaliknya). Vektor ini akan Anda
simpan di variabel `z'.

\begin{Shaded}
\begin{Highlighting}[]
\NormalTok{z <-}\StringTok{ }\KeywordTok{c}\NormalTok{(}\FloatTok{1.1}\NormalTok{, }\DecValTok{9}\NormalTok{, }\FloatTok{3.14}\NormalTok{)}
\end{Highlighting}
\end{Shaded}

Lalu, Anda tampilkan variabel `z'.

\begin{Shaded}
\begin{Highlighting}[]
\NormalTok{z}
\end{Highlighting}
\end{Shaded}

\begin{verbatim}
## [1] 1.10 9.00 3.14
\end{verbatim}

Perhatikan hasil outputnya. Tidak ada tanda koma yang memisahkan antar
nilai pada vektor tersebut.

Anda juga bisa menggabungkan vektor untuk membuat vektor baru. Misalnya,
vektor baru yang terdiri dari z, 777, dan z lagi. Tidak perlu menyimpan
di variabel, sehingga Anda bisa melihat langsung hasilnya setelah
menekan Enter.

\begin{Shaded}
\begin{Highlighting}[]
\KeywordTok{c}\NormalTok{(z, }\DecValTok{777}\NormalTok{, z)}
\end{Highlighting}
\end{Shaded}

\begin{verbatim}
## [1]   1.10   9.00   3.14 777.00   1.10   9.00   3.14
\end{verbatim}

Numerik faktor bisa digunakan dengan operator aritmatika. Ketik: z * 2 +
100, dan lihat apa yang terjadi.

\begin{Shaded}
\begin{Highlighting}[]
\NormalTok{z }\OperatorTok{*}\StringTok{ }\DecValTok{2} \OperatorTok{+}\StringTok{ }\DecValTok{100}
\end{Highlighting}
\end{Shaded}

\begin{verbatim}
## [1] 102.20 118.00 106.28
\end{verbatim}

Pertama-tama, R mengalikan 2 setiap elemen vektor z, lalu menambahkannya
dengan 100 sehingga hasilnya seperti di atas.

\subsection{Urutan Angka}\label{urutan-angka}

Cara termudah untuk membuat angka yang berurutan adalah dengan
menggunakan operator \texttt{:}. Ketik \texttt{1:20} di CONSOLE dan
tekan Enter untuk melihat apa yang akan terjadi.

\begin{Shaded}
\begin{Highlighting}[]
\DecValTok{1}\OperatorTok{:}\DecValTok{20}
\end{Highlighting}
\end{Shaded}

\begin{verbatim}
##  [1]  1  2  3  4  5  6  7  8  9 10 11 12 13 14 15 16 17 18 19 20
\end{verbatim}

Lalu bagaimana halnya dengan angka real? Misalnya, coba Anda ketik
\texttt{pi:10}.

\begin{Shaded}
\begin{Highlighting}[]
\NormalTok{pi}\OperatorTok{:}\DecValTok{10}
\end{Highlighting}
\end{Shaded}

\begin{verbatim}
## [1] 3.141593 4.141593 5.141593 6.141593 7.141593 8.141593 9.141593
\end{verbatim}

Hasilnya, angka real dimulai dari pi (3.141593), angka naik (bertambah
1) sampai dengan 9. Batas atas 10 tidak tercapai karena angka
selanjutnya (10.141593) lebih besar dari 10.

Selain \texttt{:}, Anda bisa menggunakan fungsi \texttt{seq()}.
Kelebihan fungsi ini, Anda bisa memilih opsi penambahannya. Contoh, Anda
ingin membuat urutan angka dari 1 sampai dengan 5 dengan penambahan 0.5.

\begin{Shaded}
\begin{Highlighting}[]
\KeywordTok{seq}\NormalTok{(}\DecValTok{1}\NormalTok{, }\DecValTok{5}\NormalTok{, }\DataTypeTok{by =} \FloatTok{0.5}\NormalTok{)}
\end{Highlighting}
\end{Shaded}

\begin{verbatim}
## [1] 1.0 1.5 2.0 2.5 3.0 3.5 4.0 4.5 5.0
\end{verbatim}

\subsection{Replikasi Angka}\label{replikasi-angka}

Fungsi untuk mereplika Angka adalah \texttt{rep()} yang merupakan
kepanjangan dari `replicate'. Berikut contoh mereplikasi angka 0
sebanyak 30 kali.

\begin{Shaded}
\begin{Highlighting}[]
\KeywordTok{rep}\NormalTok{(}\DecValTok{0}\NormalTok{, }\DataTypeTok{times =} \DecValTok{30}\NormalTok{)}
\end{Highlighting}
\end{Shaded}

\begin{verbatim}
##  [1] 0 0 0 0 0 0 0 0 0 0 0 0 0 0 0 0 0 0 0 0 0 0 0 0 0 0 0 0 0 0
\end{verbatim}

\subsection{Ambil Elemen Vektor}\label{ambil-elemen-vektor}

Anda bisa mengambil elemen vektor dengan menggunakan operator
\texttt{{[}{]}} dan operator urutan angka \texttt{:}, serta menggunakan
fungsi \texttt{c()}.

Misalkan, Anda ambil contoh: variabel `a' adalah vektor yang terdiri
dari nilai 2, 4, 6, dan 8.

\begin{Shaded}
\begin{Highlighting}[]
\NormalTok{a <-}\StringTok{ }\KeywordTok{c}\NormalTok{(}\DecValTok{2}\NormalTok{, }\DecValTok{4}\NormalTok{, }\DecValTok{6}\NormalTok{, }\DecValTok{8}\NormalTok{)}
\end{Highlighting}
\end{Shaded}

\begin{Shaded}
\begin{Highlighting}[]
\CommentTok{# Ambil elemen ke-1 sampai dengan yang ke-3 dari variabel a}
\NormalTok{a[}\DecValTok{1}\OperatorTok{:}\DecValTok{3}\NormalTok{]}
\end{Highlighting}
\end{Shaded}

\begin{verbatim}
## [1] 2 4 6
\end{verbatim}

\begin{Shaded}
\begin{Highlighting}[]
\CommentTok{# Ambil elemen ke-2 dari variabel a}
\NormalTok{a[}\DecValTok{2}\NormalTok{]}
\end{Highlighting}
\end{Shaded}

\begin{verbatim}
## [1] 4
\end{verbatim}

\begin{Shaded}
\begin{Highlighting}[]
\CommentTok{# Ambil elemen ke-4 dan ke-1 dari variabel a}
\NormalTok{a[}\KeywordTok{c}\NormalTok{(}\DecValTok{4}\NormalTok{, }\DecValTok{1}\NormalTok{)]}
\end{Highlighting}
\end{Shaded}

\begin{verbatim}
## [1] 8 2
\end{verbatim}

\begin{Shaded}
\begin{Highlighting}[]
\CommentTok{# Keluarkan elemen ke-3 dari variabel a}
\NormalTok{a[}\OperatorTok{-}\DecValTok{3}\NormalTok{]}
\end{Highlighting}
\end{Shaded}

\begin{verbatim}
## [1] 2 4 8
\end{verbatim}

\begin{Shaded}
\begin{Highlighting}[]
\CommentTok{# Keluarkan elemen ke-1 dan ke-3 dari variabel a}
\NormalTok{a[}\OperatorTok{-}\KeywordTok{c}\NormalTok{(}\DecValTok{1}\NormalTok{, }\DecValTok{3}\NormalTok{)]}
\end{Highlighting}
\end{Shaded}

\begin{verbatim}
## [1] 4 8
\end{verbatim}

\subsection{\texorpdfstring{\emph{Syntax} / Fungsi Penting
Lainnya}{Syntax / Fungsi Penting Lainnya}}\label{syntax-fungsi-penting-lainnya}

\begin{tabular}{l|l|l}
\hline
\_Syntax\_ & Keterangan & Contoh Penggunaan\\
\hline
? & Akses ke dokumentasi dan penjelasan fungsi & ?seq\\
\hline
getwd() & Menampilkan filepath direktori kerja & getwd()\\
\hline
ls() & Menampilkan daftar varibel & ls()\\
\hline
dir() & Menampilkan files dari sebuah direktori & dir()\\
\hline
rm(list=ls()) & Menghapus seluruh variabel & rm(list=ls())\\
\hline
length() & Mendapatkan panjang elemen sebuah vektor atau variabel & length(x)\\
\hline
\end{tabular}

\textbf{Tombol Keyboard}

Beberapa tombol keyboard berfungsi untuk mempercepat Anda dalam menulis
script di R:

\begin{itemize}
\item
  \texttt{tab} sangat sering digunakan untuk fitur \emph{auto
  completion}. Contohnya, jika Anda memiliki varibel bernama
  `aristoteles', Anda cukup ketik `a' di CONSOLE, lalu tekan tombol
  \texttt{tab}, maka R akan langsung menyempurnakan penulisan variabel
  tersebut.
\item
  \texttt{panah\ atas} dan \texttt{panah\ bawah} berfungsi untuk
  mengakses histori perintah yang sudah Anda input sebelumnya, sehingga
  Anda tidak perlu mengetik ulang.
\end{itemize}

\chapter{Kuasai Teknik Dasar Analisis
Data}\label{kuasai-teknik-dasar-analisis-data}

Analisis data merupakan subjek yang cukup luas. Anda tidak bisa belajar
semuanya sekaligus. Waktu Anda terbatas dan sangat berharga. Anda perlu
strategi khusus untuk menguasainya. Investasi terbaik Anda adalah waktu
dan pikiran. Anda harus fokus belajar satu keterampilan yang memiliki
keuntungan atau balik modal investasi yang tinggi.

Kabar baiknya, Anda sudah melakukan hal ini: Belajar R.

R merupakan tool yang memiliki keuntungan atau balik modal investasi
paling tinggi.

Jadi, lanjutkanlah belajar R.

Di bagian ini, Anda akan diarahkan untuk mengetahui sebenarnya apa saja
yang menjadi bagian dari teknik dasar analisis data.

Teknik dasar analisis data dibagi menjadi 3:

\begin{enumerate}
\def\labelenumi{\arabic{enumi}.}
\tightlist
\item
  Visualisasi Data
\item
  Manipulasi Data
\item
  Interpretasi Data
\end{enumerate}

Anda sebaiknya belajar teknik dasar analisis ini secara berurutan.

Pertama-tama pelajari visualisasi data.

Visualisasi data merupakan cara tercepat dan termudah untuk mendapatkan
\emph{insight} pada sebuah set data. Selain itu, keterampilan ini juga
bisa dijadikan sebagai alat komunikasi ke siapa saja yang membutuhkan
hasil analisis data.

\section{Visualisasi Data}\label{visualisasi-data}

Untuk belajar visualisasi data, gunakan salah satu paket R, yakni
\texttt{ggplot2}.

Huruf ``gg'' pada \texttt{ggplot2} merupakan kepanjangan dari
\emph{Grammar of Graphics}, sebuah buku yang ditulis oleh Leland
Wilkinson, prinsip yang \texttt{ggplot2} ambil. \emph{Grammar of
Graphics} menjelaskan tentang metode membagi-bagi grafik ke berbagai
elemen dan membangun setiap elemen tersebut di beberapa lapisan untuk
menampilkan representasi visual.

Singkatnya, grafik di petakan dari data ke atribut estetik, seperti
warna, bentuk, atau ukuran juga ke objek geometrik, seperti titik,
garis, ataupun batang. Selain itu, grafik juga dipetakan ke sistem
koordinat.

Sebuah grafik ggplot terdiri dari sebuah sistem koordinat dan setidaknya
sebuah objek geometrik.

Lalu, Kenapa \texttt{ggplot2}?

\texttt{ggplot2} mengajarkan Anda bagaimana berpikir tentang visualisasi
data.

Anda diajarkan tidak hanya secara \emph{syntax} saja, tetapi juga
berpikir dan memiliki konsep terkait proses memvisualisasi data.

\textbf{3 Prinsip Dasar Visualisasi Menggunakan \texttt{ggplot2}}

\begin{enumerate}
\def\labelenumi{\arabic{enumi}.}
\tightlist
\item
  Pemetaan Data ke Atribut Estetik
\item
  Lapisan
\item
  Iterasi Pemetaan dan Lapisan
\end{enumerate}

OK, sampai sini, mungkin Anda masih bingung, tidak apa. Ada baiknya
konsep ini langsung dijelaskan melalui contoh saja.

\subsection{Pemetaan Data ke Atribut
Estetik}\label{pemetaan-data-ke-atribut-estetik}

Pada contoh kali ini, Anda akan diarahkan untuk membuat grafik batang
atau \emph{barchart}. Grafik ini merupakan grafik yang paling sering
dipakai dan paling mudah dibaca. Selain itu, akan lebih memudahkan Anda
untuk memahami prinsip dan filosofi membuat grafik menggunakan
\texttt{ggplot2}.

Sebenarnya ada banyak grafik jenis lainnya seperti histogram, grafik
garis \emph{(line chart)}, grafik sebar \emph{(scatter plot)}, dan
grafik lainnya. Pembahasan di bagian ini akan difokuskan pada grafik
batang saja.

Cukup penjelasannya, coba langsung Anda inisiasi paket \texttt{ggplot2}
menggunakan fungsi \texttt{library()}.

\begin{Shaded}
\begin{Highlighting}[]
\CommentTok{# Inisiasi paket ggplot}
\KeywordTok{library}\NormalTok{(ggplot2)}
\end{Highlighting}
\end{Shaded}

Setiap Anda ingin menggunakan paket R yang cukup spesifik, Anda harus
menginisasi paket tersebut di awal. Hal ini dilakukan agar Anda bisa
mengakses seluruh fungsi yang ada di paket tersebut.

Anda akan menggunakan data \emph{built in} dari \texttt{ggplot2}, yakni
\texttt{mpg}. Untuk melihat data \texttt{mpg}, ketik \texttt{mpg} di
CONSOLE dan tekan Enter.

\begin{Shaded}
\begin{Highlighting}[]
\NormalTok{mpg}
\end{Highlighting}
\end{Shaded}

\begin{verbatim}
## # A tibble: 234 x 11
##    manufacturer model displ  year   cyl trans drv     cty   hwy fl    cla~
##    <chr>        <chr> <dbl> <int> <int> <chr> <chr> <int> <int> <chr> <ch>
##  1 audi         a4      1.8  1999     4 auto~ f        18    29 p     com~
##  2 audi         a4      1.8  1999     4 manu~ f        21    29 p     com~
##  3 audi         a4      2    2008     4 manu~ f        20    31 p     com~
##  4 audi         a4      2    2008     4 auto~ f        21    30 p     com~
##  5 audi         a4      2.8  1999     6 auto~ f        16    26 p     com~
##  6 audi         a4      2.8  1999     6 manu~ f        18    26 p     com~
##  7 audi         a4      3.1  2008     6 auto~ f        18    27 p     com~
##  8 audi         a4 q~   1.8  1999     4 manu~ 4        18    26 p     com~
##  9 audi         a4 q~   1.8  1999     4 auto~ 4        16    25 p     com~
## 10 audi         a4 q~   2    2008     4 manu~ 4        20    28 p     com~
## # ... with 224 more rows
\end{verbatim}

Seperti yang Anda lihat, \texttt{mpg}, yang merupakan kepanjangan dari
\emph{miles per gallon} adalah sebuah set data tentang efisiensi
penggunaan bensin 38 mobil terpopuler di Amerika Serikat.

Anda bisa memahami lebih lanjut terkait detail set data ini melalui
perintah \texttt{?mpg}.

Untuk melihat berapa jumlah baris dan kolomnya, Anda bisa menggunakan
fungsi \texttt{dim}.

\begin{Shaded}
\begin{Highlighting}[]
\KeywordTok{dim}\NormalTok{(mpg)}
\end{Highlighting}
\end{Shaded}

\begin{verbatim}
## [1] 234  11
\end{verbatim}

Set data ini memiliki 234 baris dan 11 kolom atau 234 observasi dan 11
variabel. Jika Anda lihat data ini lebih teliti lagi, 234 juga merupakan
jumlah total keseluruhan mobil yang menjadi objek observasi.

Sebelum melangkah lebih jauh ke teknis pembuatan grafik, Anda akan
diarahkan untuk mengetahui \textbf{perbandingan jumlah mobil antar
`manufacturer'} sebagai tujuan dari proses ini.

OK, untuk mulai memetakan variabel ke grafik, gunakan fungsi
\texttt{ggplot()}.

\begin{Shaded}
\begin{Highlighting}[]
\CommentTok{# Pemetaan variabel `manufacturer` ke sistem koordinat x}
\KeywordTok{ggplot}\NormalTok{(}\DataTypeTok{data =}\NormalTok{ mpg, }\KeywordTok{aes}\NormalTok{(}\DataTypeTok{x =}\NormalTok{ manufacturer))}
\end{Highlighting}
\end{Shaded}

\includegraphics{ebook_files/figure-latex/unnamed-chunk-27-1.pdf}

Fungsi \texttt{aes()} di atas adalah memetakan atribut variabel
`manufacturer' ke sistem koordinat x. Anda bisa lihat, di garis
horizontal x sudah ada nama-nama pembuat mobil yang merupakan hasil
pemetaan atribut estetik.

Sampai sini Anda pasti berpikir,

\emph{``Kenapa kok jelek banget grafiknya?''}

\emph{``Mana grafik batangnya?''}

\emph{``Tulisannya tumpang tindih, nggak beraturan.''}

Tenang, grafik tersebut masih tahap lapisan pertama dan iterasi pertama
juga prosesnya masih langkah demi langkah sehingga Anda tidak bingung
dan lebih mudah mengerti konsep membuat grafik \texttt{ggplot2}.

\subsection{Lapisan}\label{lapisan}

Lapisan atau \emph{layer} merupakan prinsip dasar visualisasi kedua
menggunakan \texttt{ggplot2}. Filosofinya adalah Anda membuat grafik
yang terdiri dari beberapa lapisan sehingga hasilnya merepresentasikan
visualisasi yang Anda inginkan. Lapisan ini bisa berupa sistem
koordinat, grafik batang, judul grafik, legenda, ataupun elemen lainnya.

Langkah selanjutnya setelah pemetaan variabel, Anda akan menambahkan
elemen grafik batang pada lapisan berikutnya dengan menggunakan operator
\texttt{+} dan fungsi \texttt{geom\_bar()}.

\begin{Shaded}
\begin{Highlighting}[]
\CommentTok{# Penambahan elemen grafik batang pada lapisan kedua}
\KeywordTok{ggplot}\NormalTok{(}\DataTypeTok{data =}\NormalTok{ mpg, }\KeywordTok{aes}\NormalTok{(}\DataTypeTok{x =}\NormalTok{ manufacturer)) }\OperatorTok{+}\StringTok{ }
\StringTok{  }\KeywordTok{geom_bar}\NormalTok{()}
\end{Highlighting}
\end{Shaded}

\includegraphics{ebook_files/figure-latex/unnamed-chunk-28-1.pdf}

Nah, sudah mulai terlihat grafiknya.

Sampai sini, sudah ada 2 lapisan grafik. Lapisan pertama diisi oleh
elemen sistem koordinat, lapisan kedua diisi oleh grafik batang.

\subsection{Iterasi Pemetaan dan
Lapisan}\label{iterasi-pemetaan-dan-lapisan}

Prinsip dasar ketiga ini merupakan perulangan dari prinsip satu dan
prinsip kedua. Anda terus melakukan iterasi dan modifikasi atribut
estetik serta penambahan lapisan atau elemen sehingga sesuai dengan
representasi visual yang ingin Anda capai.

Untuk mempermudah pemahaman, Anda akan diarahkan langsung melalui contoh
penulisan kode yang disertai dengan penjelasan dari baris komentar
(baris komentar diawali dengan tanda \texttt{\#}).

\begin{Shaded}
\begin{Highlighting}[]
\CommentTok{# Merubah warna grafik batang menjadi merah gelap}
\CommentTok{# dengan menggunakan argumen "fill" pada fungsi geom_bar()}
\KeywordTok{ggplot}\NormalTok{(}\DataTypeTok{data =}\NormalTok{ mpg, }\KeywordTok{aes}\NormalTok{(}\DataTypeTok{x =}\NormalTok{ manufacturer)) }\OperatorTok{+}\StringTok{ }
\StringTok{  }\KeywordTok{geom_bar}\NormalTok{(}\DataTypeTok{fill =} \StringTok{"#8B0000"}\NormalTok{)}
\end{Highlighting}
\end{Shaded}

\includegraphics{ebook_files/figure-latex/unnamed-chunk-29-1.pdf}

\begin{Shaded}
\begin{Highlighting}[]
\CommentTok{# Merubah label keterangan sistem koordinat x dan y}
\CommentTok{# dengan menggunakan fungsi labs()}
\KeywordTok{ggplot}\NormalTok{(}\DataTypeTok{data =}\NormalTok{ mpg, }\KeywordTok{aes}\NormalTok{(}\DataTypeTok{x =}\NormalTok{ manufacturer)) }\OperatorTok{+}\StringTok{ }
\StringTok{  }\KeywordTok{geom_bar}\NormalTok{(}\DataTypeTok{fill =} \StringTok{"#8B0000"}\NormalTok{) }\OperatorTok{+}
\StringTok{  }\KeywordTok{labs}\NormalTok{(}\DataTypeTok{x =} \StringTok{"Produsen"}\NormalTok{, }\DataTypeTok{y =} \StringTok{"Jumlah Mobil"}\NormalTok{)}
\end{Highlighting}
\end{Shaded}

\includegraphics{ebook_files/figure-latex/unnamed-chunk-30-1.pdf}

\begin{Shaded}
\begin{Highlighting}[]
\CommentTok{# Menambah judul dan subjudul}
\CommentTok{# dengan menggunakan fungsi ggtitle()}
\KeywordTok{ggplot}\NormalTok{(}\DataTypeTok{data =}\NormalTok{ mpg, }\KeywordTok{aes}\NormalTok{(}\DataTypeTok{x =}\NormalTok{ manufacturer)) }\OperatorTok{+}\StringTok{ }
\StringTok{  }\KeywordTok{geom_bar}\NormalTok{(}\DataTypeTok{fill =} \StringTok{"#8B0000"}\NormalTok{) }\OperatorTok{+}
\StringTok{  }\KeywordTok{labs}\NormalTok{(}\DataTypeTok{x =} \StringTok{"Produsen"}\NormalTok{, }\DataTypeTok{y =} \StringTok{"Jumlah Mobil"}\NormalTok{) }\OperatorTok{+}
\StringTok{  }\KeywordTok{ggtitle}\NormalTok{(}\StringTok{"Distribusi Jumlah Mobil Berdasarkan Produsen"}\NormalTok{, }
          \DataTypeTok{subtitle =} \StringTok{"Tahun Pembuatan 1999 dan 2008"}\NormalTok{)}
\end{Highlighting}
\end{Shaded}

\includegraphics{ebook_files/figure-latex/unnamed-chunk-31-1.pdf}

Begitulah bagaimana iterasi dilakukan terus menerus sehingga
representasi grafik sesuai dengan yang Anda inginkan.

Dan pada akhirnya, iterasi final, menukar sistem koordinat sehingga
tulisan terlihat lebih rapih dan mudah dibaca.

\begin{Shaded}
\begin{Highlighting}[]
\CommentTok{# Menukar sistem koordinat sehingga lebih mudah dibaca}
\CommentTok{# menggunakan fungsi coord_flip()}
\KeywordTok{ggplot}\NormalTok{(}\DataTypeTok{data =}\NormalTok{ mpg, }\KeywordTok{aes}\NormalTok{(}\DataTypeTok{x =}\NormalTok{ manufacturer)) }\OperatorTok{+}\StringTok{ }
\StringTok{  }\KeywordTok{geom_bar}\NormalTok{(}\DataTypeTok{fill =} \StringTok{"#8B0000"}\NormalTok{) }\OperatorTok{+}
\StringTok{  }\KeywordTok{labs}\NormalTok{(}\DataTypeTok{x =} \StringTok{"Produsen"}\NormalTok{, }\DataTypeTok{y =} \StringTok{"Jumlah Mobil"}\NormalTok{) }\OperatorTok{+}
\StringTok{  }\KeywordTok{ggtitle}\NormalTok{(}\StringTok{"Distribusi Jumlah Mobil Berdasarkan Produsen"}\NormalTok{, }
          \DataTypeTok{subtitle =} \StringTok{"Tahun Pembuatan 1999 dan 2008"}\NormalTok{) }\OperatorTok{+}
\StringTok{  }\KeywordTok{coord_flip}\NormalTok{()}
\end{Highlighting}
\end{Shaded}

\includegraphics{ebook_files/figure-latex/unnamed-chunk-32-1.pdf}

Sekali Anda mengerti tentang prinsip dasar ini, pemahaman Anda terhadap
data visualisasi akan berubah. Mungkin bagi Anda yang dulu terbiasa
menggunakan grafik dari Excel akan memerlukan waktu lebih untuk
menangkap ini semua. Bagaimana tidak, Excel memberikan kemudahan yang
membuat Anda tidak perlu memikirkan proses bagaimana grafik dibuat. Anda
hanya cukup \emph{point, drag, and click}, sedikit konfigurasi, Excel
akan membereskan semuanya.

Pelajari \texttt{ggplot2}, Anda akan paham bagaimana proses grafik
dibuat tahap demi tahap. Memang membutuhkan \emph{effort} lebih, tetapi
\emph{worth it}. Seiring Anda mengerti cara menggunakan
\texttt{ggplot2}, pemahaman prinsip visualisasi data Anda akan
meningkat.

\section{Manipulasi Data}\label{manipulasi-data}

Manipulasi data berperan ketika representasi visualisasi yang Anda buat
tidak cukup untuk memunculkan \emph{insight}. Faktanya, data yang Anda
peroleh hampir pasti perlu diolah terlebih dahulu. Menurut penelitian,
seorang analis data kebanyakan mengerahkan usahanya berkutat di fase
manipulasi atau mengolah data ini sekitar 60\% - 70\% dari keseluruhan
proses analisis. Jadi bisa dikatakan, keterampilan manipulasi data cukup
penting dan pastinya juga perlu Anda kuasai.

Manipulasi data adalah segala proses yang dilakukan pada sebuah set data
dengan tujuan mempermudah melakukan analisis data dan mendapatkan
\emph{insight}. Proses tersebut diantaranya mengubah nama kolom atau
variabel, mengurutkan, menyaring, sampai dengan meringkas data.

Konsep atau teknik manipulasi data yang akan dikupas diantaranya sebagai
berikut:

\begin{enumerate}
\def\labelenumi{\arabic{enumi}.}
\tightlist
\item
  Melihat dan Mengetahui Set Data
\item
  Mengubah Variabel, Kolom, Isi Data
\item
  Operator Pipa (\texttt{\%\textgreater{}\%})
\item
  Memilih Variabel Data
\item
  Menyaring Baris Data
\item
  Mengurutkan Data
\item
  Mengelompokkan Data
\item
  Meringkas Data
\end{enumerate}

Pada bahasan kali ini, Anda akan diarahkan untuk menggunakan kembali set
data \texttt{mpg} seperti sebelumnya ditambah satu paket manipulasi data
yang sangat \emph{powerful}, \texttt{dplyr}.

Pertama-tama, inisiasi paket \texttt{dplyr}.

\begin{Shaded}
\begin{Highlighting}[]
\CommentTok{# Inisiasi paket dplyr}
\KeywordTok{library}\NormalTok{(dplyr)}
\end{Highlighting}
\end{Shaded}

\subsection{Melihat dan Mengetahui Set
Data}\label{melihat-dan-mengetahui-set-data}

Setiap Anda mendapatkan sebuah data set untuk diolah, hal yang pertama
harus Anda lakukan adalah mengetahui isi dan tentang apa data tersebut.
Apa format datanya? Berapa dimensinya? Apa saja variabelnya? Bagaimana
tipe variabelnya? Apakah ada data yang hilang atau tidak terisi? Dan
lain sebagainya.

Yang pertama adalah fungsi \texttt{class()} untuk menampilkan tipe data.

\begin{Shaded}
\begin{Highlighting}[]
\CommentTok{# Menampilkan tipe data}
\KeywordTok{class}\NormalTok{(mpg)}
\end{Highlighting}
\end{Shaded}

\begin{verbatim}
## [1] "tbl_df"     "tbl"        "data.frame"
\end{verbatim}

Apa itu ``tbl\_df'' ``tbl''?

``tbl\_df'' = \emph{table data frame}

``tbl'' = \emph{table}

Intinya, ``tbl\_df'' ``tbl'' ``data.frame'' merupakan tipe data dari
\texttt{mpg} yang berupa tabel atau data frame. Tabel atau data frame
adalah tipe data berbentuk ``kotak'' yang terdiri dari baris dan kolom.

Data yang mudah diolah atau dimanipulasi biasanya bertipe tabel atau
data frame ini. Selain tabel, ada juga yang berbentuk vektor atau yang
hanya terdiri dari satu baris saja atau satu kolom saja.

Lalu, bagaimana cara mengetahui jumlah baris dan kolom data
\texttt{mpg}? Ya, Anda benar, salah satunya menggunakan fungsi
\texttt{dim()}.

Anda juga bisa menggunakan fungsi \texttt{nrow()} dan \texttt{ncol()}.

\begin{Shaded}
\begin{Highlighting}[]
\CommentTok{# Menampilkan jumlah baris}
\KeywordTok{nrow}\NormalTok{(mpg)}
\end{Highlighting}
\end{Shaded}

\begin{verbatim}
## [1] 234
\end{verbatim}

\begin{Shaded}
\begin{Highlighting}[]
\CommentTok{# Menampilkan jumlah kolom atau variabel}
\KeywordTok{ncol}\NormalTok{(mpg)}
\end{Highlighting}
\end{Shaded}

\begin{verbatim}
## [1] 11
\end{verbatim}

Dari fungsi di atas, Anda sudah mengetahui jumlah kolom atau variabel.
Selanjutnya, Anda gali lebih jauh untuk mengetahui nama-nama kolom atau
variabelnya. Fungsi \texttt{names()} bisa membantu Anda.

\begin{Shaded}
\begin{Highlighting}[]
\CommentTok{# Menampilkan seluruh variabel atau nama kolom}
\KeywordTok{names}\NormalTok{(mpg)}
\end{Highlighting}
\end{Shaded}

\begin{verbatim}
##  [1] "manufacturer" "model"        "displ"        "year"        
##  [5] "cyl"          "trans"        "drv"          "cty"         
##  [9] "hwy"          "fl"           "class"
\end{verbatim}

Anda sudah tahu variabel dari set data \texttt{mpg}, akan lebih lengkap
lagi jika Anda tahu isi di dalamnya. Gunakan fungsi \texttt{head()} dan
\texttt{tail()} untuk mengintip isi bagian awal dan akhir dari sebuah
set data.

\begin{Shaded}
\begin{Highlighting}[]
\CommentTok{# Menampilkan 7 baris awal dari data}
\KeywordTok{head}\NormalTok{(mpg, }\DecValTok{7}\NormalTok{)}
\end{Highlighting}
\end{Shaded}

\begin{verbatim}
## # A tibble: 7 x 11
##   manufacturer model displ  year   cyl trans drv     cty   hwy fl    class
##   <chr>        <chr> <dbl> <int> <int> <chr> <chr> <int> <int> <chr> <chr>
## 1 audi         a4      1.8  1999     4 auto~ f        18    29 p     comp~
## 2 audi         a4      1.8  1999     4 manu~ f        21    29 p     comp~
## 3 audi         a4      2    2008     4 manu~ f        20    31 p     comp~
## 4 audi         a4      2    2008     4 auto~ f        21    30 p     comp~
## 5 audi         a4      2.8  1999     6 auto~ f        16    26 p     comp~
## 6 audi         a4      2.8  1999     6 manu~ f        18    26 p     comp~
## 7 audi         a4      3.1  2008     6 auto~ f        18    27 p     comp~
\end{verbatim}

\begin{Shaded}
\begin{Highlighting}[]
\CommentTok{# Menampilkan 8 baris akhir dari data}
\KeywordTok{tail}\NormalTok{(mpg, }\DecValTok{8}\NormalTok{)}
\end{Highlighting}
\end{Shaded}

\begin{verbatim}
## # A tibble: 8 x 11
##   manufacturer model displ  year   cyl trans drv     cty   hwy fl    class
##   <chr>        <chr> <dbl> <int> <int> <chr> <chr> <int> <int> <chr> <chr>
## 1 volkswagen   new ~   2.5  2008     5 auto~ f        20    29 r     subc~
## 2 volkswagen   pass~   1.8  1999     4 manu~ f        21    29 p     mids~
## 3 volkswagen   pass~   1.8  1999     4 auto~ f        18    29 p     mids~
## 4 volkswagen   pass~   2    2008     4 auto~ f        19    28 p     mids~
## 5 volkswagen   pass~   2    2008     4 manu~ f        21    29 p     mids~
## 6 volkswagen   pass~   2.8  1999     6 auto~ f        16    26 p     mids~
## 7 volkswagen   pass~   2.8  1999     6 manu~ f        18    26 p     mids~
## 8 volkswagen   pass~   3.6  2008     6 auto~ f        17    26 p     mids~
\end{verbatim}

Selain \texttt{head()} dan \texttt{tail()}, untuk melihat isi data
berformat \emph{spreadsheet}, Anda bisa menggunakan fungsi
\texttt{View()}.

\begin{Shaded}
\begin{Highlighting}[]
\CommentTok{# Menampilkan isi data berformat spreadsheet}
\KeywordTok{View}\NormalTok{(mpg)}
\end{Highlighting}
\end{Shaded}

Hasilnya akan muncul di bagian jendela SCRIPT RStudio.

\textbf{RINGKASAN DATA}

Anda sudah mengetahui beberapa cara melihat dan menggali informasi dari
sebuah set data. Ada satu cara untuk menampilkan ringkasan umum sehingga
Anda bisa langsung melihat dan menangkap isi data secara \emph{big
picture}.

\begin{Shaded}
\begin{Highlighting}[]
\CommentTok{# Menampilkan ringkasan umum dari data}
\KeywordTok{summary}\NormalTok{(mpg)}
\end{Highlighting}
\end{Shaded}

\begin{verbatim}
##  manufacturer          model               displ            year     
##  Length:234         Length:234         Min.   :1.600   Min.   :1999  
##  Class :character   Class :character   1st Qu.:2.400   1st Qu.:1999  
##  Mode  :character   Mode  :character   Median :3.300   Median :2004  
##                                        Mean   :3.472   Mean   :2004  
##                                        3rd Qu.:4.600   3rd Qu.:2008  
##                                        Max.   :7.000   Max.   :2008  
##       cyl           trans               drv                 cty       
##  Min.   :4.000   Length:234         Length:234         Min.   : 9.00  
##  1st Qu.:4.000   Class :character   Class :character   1st Qu.:14.00  
##  Median :6.000   Mode  :character   Mode  :character   Median :17.00  
##  Mean   :5.889                                         Mean   :16.86  
##  3rd Qu.:8.000                                         3rd Qu.:19.00  
##  Max.   :8.000                                         Max.   :35.00  
##       hwy             fl               class          
##  Min.   :12.00   Length:234         Length:234        
##  1st Qu.:18.00   Class :character   Class :character  
##  Median :24.00   Mode  :character   Mode  :character  
##  Mean   :23.44                                        
##  3rd Qu.:27.00                                        
##  Max.   :44.00
\end{verbatim}

Fungsi \texttt{summary()} memberikan Anda ringkasan umum dari sebuah set
data. Jika variabelnya bertipe numerik, akan ditampilkan ringkasan
seperti nilai minimal, maksimal, rata-rata, median, kuartil pertama, dan
kuartil ketiga.

Selain fungsi \texttt{summary()}, ada fungsi lain yang lebih ringkas,
yaitu \texttt{str()} dan \texttt{glimpse()}.

\begin{Shaded}
\begin{Highlighting}[]
\CommentTok{# Menampilkan struktur dan isi data}
\KeywordTok{str}\NormalTok{(mpg)}
\end{Highlighting}
\end{Shaded}

\begin{verbatim}
## Classes 'tbl_df', 'tbl' and 'data.frame':    234 obs. of  11 variables:
##  $ manufacturer: chr  "audi" "audi" "audi" "audi" ...
##  $ model       : chr  "a4" "a4" "a4" "a4" ...
##  $ displ       : num  1.8 1.8 2 2 2.8 2.8 3.1 1.8 1.8 2 ...
##  $ year        : int  1999 1999 2008 2008 1999 1999 2008 1999 1999 2008 ...
##  $ cyl         : int  4 4 4 4 6 6 6 4 4 4 ...
##  $ trans       : chr  "auto(l5)" "manual(m5)" "manual(m6)" "auto(av)" ...
##  $ drv         : chr  "f" "f" "f" "f" ...
##  $ cty         : int  18 21 20 21 16 18 18 18 16 20 ...
##  $ hwy         : int  29 29 31 30 26 26 27 26 25 28 ...
##  $ fl          : chr  "p" "p" "p" "p" ...
##  $ class       : chr  "compact" "compact" "compact" "compact" ...
\end{verbatim}

\begin{Shaded}
\begin{Highlighting}[]
\CommentTok{# Menampilkan struktur dan isi data}
\NormalTok{dplyr}\OperatorTok{::}\KeywordTok{glimpse}\NormalTok{(mpg)}
\end{Highlighting}
\end{Shaded}

\begin{verbatim}
## Observations: 234
## Variables: 11
## $ manufacturer <chr> "audi", "audi", "audi", "audi", "audi", "audi", "...
## $ model        <chr> "a4", "a4", "a4", "a4", "a4", "a4", "a4", "a4 qua...
## $ displ        <dbl> 1.8, 1.8, 2.0, 2.0, 2.8, 2.8, 3.1, 1.8, 1.8, 2.0,...
## $ year         <int> 1999, 1999, 2008, 2008, 1999, 1999, 2008, 1999, 1...
## $ cyl          <int> 4, 4, 4, 4, 6, 6, 6, 4, 4, 4, 4, 6, 6, 6, 6, 6, 6...
## $ trans        <chr> "auto(l5)", "manual(m5)", "manual(m6)", "auto(av)...
## $ drv          <chr> "f", "f", "f", "f", "f", "f", "f", "4", "4", "4",...
## $ cty          <int> 18, 21, 20, 21, 16, 18, 18, 18, 16, 20, 19, 15, 1...
## $ hwy          <int> 29, 29, 31, 30, 26, 26, 27, 26, 25, 28, 27, 25, 2...
## $ fl           <chr> "p", "p", "p", "p", "p", "p", "p", "p", "p", "p",...
## $ class        <chr> "compact", "compact", "compact", "compact", "comp...
\end{verbatim}

Seperti yang Anda lihat, fungsi \texttt{glimpse()} sepertinya lebih
rapih dalam tampilan.

Kedua fungsi ini menampilkan cukup informasi dalam sebuah tampilan yang
padat dan jelas mulai dari berapa jumlah variabel, jumlah observasi,
nama variabel, tipe variabel, sampai dengan isi dari sebagian data.

Ketika Anda mendapatkan sebuah set data untuk diolah, tentu fungsi ini
tidak akan Anda lewatkan begitu saja.

Yang perlu Anda cermati, set data \texttt{mpg} ini tipe variabelnya
masih kurang tepat. Misalnya, variabel \texttt{year} atau tahun
pembuatan mobil. Di data hanya ada tahun 1999 dan 2008 saja. Seharusnya
ini termasuk ke dalam tipe variabel faktor atau kategorikal, bukan
numerik.

Di bagian berikutnya, akan dibahas bagaimana cara mengubah variabel.

\subsection{Mengubah Variabel, Kolom, Isi
Data}\label{mengubah-variabel-kolom-isi-data}

Sebelum melangkah lebih jauh, Anda harus mengetahui beberapa tipe
variabel di pemrograman R.

\textbf{TIPE VARIABEL}

Yang sering digunakan ada 6 tipe:

\begin{itemize}
\tightlist
\item
  Karakter: Huruf, kata, atau kalimat seperti \texttt{b}, \texttt{apel},
  \texttt{mantap\ jiwa}.
\item
  Integer: Bilangan bulat, contoh \texttt{1}, \texttt{4}, \texttt{7}.
\item
  Double: Bilangan tidak bulat atau real, contoh \texttt{2.3},
  \texttt{3.5}.
\item
  Logikal: Benar atau salah (\texttt{TRUE} / \texttt{FALSE}).
\item
  Kompleks: \texttt{2+4i}, angka kompleks gabungan real dan imaginer.
\item
  Faktor: Kategorikal, seperti \texttt{laki-laki}, \texttt{perempuan}
  (jenis kelamin).
\end{itemize}

Jika Anda lihat kembali set data \texttt{mpg} menggunakan \texttt{str()}
atau \texttt{glimpse()}, maka tipe masing-masing variabelnya sebagai
berikut:

\begin{itemize}
\tightlist
\item
  manufacturer: karakter
\item
  model: karakter
\item
  displ: double
\item
  year: integer
\item
  cyl: integer
\item
  trans: karakter
\item
  drv: karakter
\item
  cty: integer
\item
  hwy: integer
\item
  fl: karakter
\item
  class: karakter
\end{itemize}

Beberapa tipe variabel di atas ada yang tidak sesuai sehingga perlu
diubah. Agar Anda lebih mengerti dan lebih jelas tentang variabel dan
isinya serta tipe variabel yang seharusnya digunakan, maka coba Anda
simak tabel berikut.

\begin{tabular}{l|l|l|l}
\hline
Variabel & Penjelasan & Tipe Variabel & Revisi Tipe Variabel\\
\hline
manufacturer & Produsen pembuat mobil & karakter & faktor\\
\hline
model & Nama model mobil & karakter & karakter\\
\hline
displ & \_Engine displacement\_ atau kapasitas mesin dalam liter & double & double\\
\hline
year & Tahun pembuatan mobil & integer & faktor\\
\hline
cyl & Jumlah silinder mesin & integer & faktor\\
\hline
trans & Tipe transmisi, auto atau manual & karakter & faktor\\
\hline
drv & Roda penggerak, f = roda depan, r = roda belakang, 4 = 4 roda & karakter & faktor\\
\hline
cty & Konsumsi bensin dalam kota (mil per galon) & integer & integer\\
\hline
hwy & Konsumsi bensin di jalan tol (mil per galon) & integer & integer\\
\hline
fl & Jenis bensin, e = ethanol, d = diesel, r = reguler, p = premium, c = natural gas & karakter & faktor\\
\hline
class & Tipe atau kelas mobil & karakter & faktor\\
\hline
\end{tabular}

Beberapa variabel seperti manufacturer, year, cyl, trans, drv, fl, dan
class merupakan tipe variabel berupa kategori sehingga perlu diubah
menjadi tipe faktor. Di set data original, variabel-variabel tersebut
dikenali sebagai tipe karakter atau integer.

\textbf{MENGAKSES VARIABEL ATAU KOLOM DALAM SET DATA}

Untuk mengubah variabel, Anda harus tahu cara mengakses salah satu
variabel atau kolom beserta isi dari sebuah set data. Operator yang
sering digunakan untuk mengakses variabel atau kolom adalah \texttt{\$}.
Ketik nama set data diikuti operator \texttt{\$} lalu nama variabel atau
kolom. Semua diketik tanpa spasi.

\begin{Shaded}
\begin{Highlighting}[]
\CommentTok{# Mengakses variabel atau kolom drv}
\NormalTok{mpg}\OperatorTok{$}\NormalTok{drv}
\end{Highlighting}
\end{Shaded}

\begin{verbatim}
##   [1] "f" "f" "f" "f" "f" "f" "f" "4" "4" "4" "4" "4" "4" "4" "4" "4" "4"
##  [18] "4" "r" "r" "r" "r" "r" "r" "r" "r" "r" "r" "4" "4" "4" "4" "f" "f"
##  [35] "f" "f" "f" "f" "f" "f" "f" "f" "f" "f" "f" "f" "f" "f" "4" "4" "4"
##  [52] "4" "4" "4" "4" "4" "4" "4" "4" "4" "4" "4" "4" "4" "4" "4" "4" "4"
##  [69] "4" "4" "4" "4" "4" "4" "r" "r" "r" "4" "4" "4" "4" "4" "4" "4" "4"
##  [86] "4" "4" "4" "4" "4" "r" "r" "r" "r" "r" "r" "r" "r" "r" "f" "f" "f"
## [103] "f" "f" "f" "f" "f" "f" "f" "f" "f" "f" "f" "f" "f" "f" "f" "f" "f"
## [120] "f" "f" "f" "4" "4" "4" "4" "4" "4" "4" "4" "4" "4" "4" "4" "r" "r"
## [137] "r" "4" "4" "4" "4" "f" "f" "f" "f" "f" "f" "f" "f" "f" "4" "4" "4"
## [154] "4" "f" "f" "f" "f" "f" "4" "4" "4" "4" "4" "4" "4" "4" "4" "4" "4"
## [171] "4" "4" "4" "4" "4" "4" "4" "4" "4" "f" "f" "f" "f" "f" "f" "f" "f"
## [188] "f" "f" "f" "f" "f" "f" "f" "f" "f" "f" "f" "4" "4" "4" "4" "4" "4"
## [205] "4" "4" "4" "f" "f" "f" "f" "f" "f" "f" "f" "f" "f" "f" "f" "f" "f"
## [222] "f" "f" "f" "f" "f" "f" "f" "f" "f" "f" "f" "f" "f"
\end{verbatim}

Seluruh 234 isi dari variabel drv ditampilkan. Anda juga bisa mencoba
variabel yang lainnya dengan cara yang sama.

\textbf{MENGUBAH TIPE VARIABEL}

Sebelum melakukan manipulasi data, pertama-tama pastikan Anda menyalin
set data asli ke set data yang baru.

\begin{Shaded}
\begin{Highlighting}[]
\CommentTok{# Menyalin set data mpg ke set data baru bernama df.mobil}
\CommentTok{# df merupakan singkatan dari data frame}
\NormalTok{df.mobil <-}\StringTok{ }\NormalTok{mpg}
\end{Highlighting}
\end{Shaded}

Mulai sekarang dan selanjutnya, segala manipulasi data dilakukan di set
data \texttt{df.mobil}.

Untuk mengubah tipe variabel ke tipe faktor atau kategorikal digunakan
fungsi \texttt{as.factor()}. Fungsi pengubah tipe variabel lainnya:

\begin{itemize}
\tightlist
\item
  \texttt{as.character()} mengubah ke tipe karakter
\item
  \texttt{as.integer()} mengubah ke tipe integer
\item
  \texttt{as.double()} mengubah ke tipe double
\item
  \texttt{as.logical()} mengubah ke tipe logikal
\item
  \texttt{as.complex()} mengubah ke tipe kompleks
\end{itemize}

\begin{Shaded}
\begin{Highlighting}[]
\CommentTok{# Mengubah tipe variabel manufacturer, year, cyl, trans, drv, fl, dan class}
\CommentTok{# ke tipe faktor}
\NormalTok{df.mobil}\OperatorTok{$}\NormalTok{manufacturer <-}\StringTok{ }\KeywordTok{as.factor}\NormalTok{(df.mobil}\OperatorTok{$}\NormalTok{manufacturer)}
\NormalTok{df.mobil}\OperatorTok{$}\NormalTok{year <-}\StringTok{ }\KeywordTok{as.factor}\NormalTok{(df.mobil}\OperatorTok{$}\NormalTok{year)}
\NormalTok{df.mobil}\OperatorTok{$}\NormalTok{cyl <-}\StringTok{ }\KeywordTok{as.factor}\NormalTok{(df.mobil}\OperatorTok{$}\NormalTok{cyl)}
\NormalTok{df.mobil}\OperatorTok{$}\NormalTok{trans <-}\StringTok{ }\KeywordTok{as.factor}\NormalTok{(df.mobil}\OperatorTok{$}\NormalTok{trans)}
\NormalTok{df.mobil}\OperatorTok{$}\NormalTok{drv <-}\StringTok{ }\KeywordTok{as.factor}\NormalTok{(df.mobil}\OperatorTok{$}\NormalTok{drv)}
\NormalTok{df.mobil}\OperatorTok{$}\NormalTok{fl <-}\StringTok{ }\KeywordTok{as.factor}\NormalTok{(df.mobil}\OperatorTok{$}\NormalTok{fl)}
\NormalTok{df.mobil}\OperatorTok{$}\NormalTok{class <-}\StringTok{ }\KeywordTok{as.factor}\NormalTok{(df.mobil}\OperatorTok{$}\NormalTok{class)}
\end{Highlighting}
\end{Shaded}

Hasilnya bisa dicek melalui fungsi \texttt{str()} atau glimpse.

\begin{Shaded}
\begin{Highlighting}[]
\CommentTok{# Menampilkan struktur set data df.mobil}
\KeywordTok{str}\NormalTok{(df.mobil)}
\end{Highlighting}
\end{Shaded}

\begin{verbatim}
## Classes 'tbl_df', 'tbl' and 'data.frame':    234 obs. of  11 variables:
##  $ manufacturer: Factor w/ 15 levels "audi","chevrolet",..: 1 1 1 1 1 1 1 1 1 1 ...
##  $ model       : chr  "a4" "a4" "a4" "a4" ...
##  $ displ       : num  1.8 1.8 2 2 2.8 2.8 3.1 1.8 1.8 2 ...
##  $ year        : Factor w/ 2 levels "1999","2008": 1 1 2 2 1 1 2 1 1 2 ...
##  $ cyl         : Factor w/ 4 levels "4","5","6","8": 1 1 1 1 3 3 3 1 1 1 ...
##  $ trans       : Factor w/ 10 levels "auto(av)","auto(l3)",..: 4 9 10 1 4 9 1 9 4 10 ...
##  $ drv         : Factor w/ 3 levels "4","f","r": 2 2 2 2 2 2 2 1 1 1 ...
##  $ cty         : int  18 21 20 21 16 18 18 18 16 20 ...
##  $ hwy         : int  29 29 31 30 26 26 27 26 25 28 ...
##  $ fl          : Factor w/ 5 levels "c","d","e","p",..: 4 4 4 4 4 4 4 4 4 4 ...
##  $ class       : Factor w/ 7 levels "2seater","compact",..: 2 2 2 2 2 2 2 2 2 2 ...
\end{verbatim}

Variabel-variabel terkait yang sebelumnya bertipe integer atau karakter
telah diubah menjadi faktor.

Perubahan lain yang bisa dilihat melalui fungsi \texttt{summary()}.

\begin{Shaded}
\begin{Highlighting}[]
\CommentTok{# Menampilkan ringkasan umum dari set data df.mobil}
\KeywordTok{summary}\NormalTok{(df.mobil)}
\end{Highlighting}
\end{Shaded}

\begin{verbatim}
##      manufacturer    model               displ         year     cyl   
##  dodge     :37    Length:234         Min.   :1.600   1999:117   4:81  
##  toyota    :34    Class :character   1st Qu.:2.400   2008:117   5: 4  
##  volkswagen:27    Mode  :character   Median :3.300              6:79  
##  ford      :25                       Mean   :3.472              8:70  
##  chevrolet :19                       3rd Qu.:4.600                    
##  audi      :18                       Max.   :7.000                    
##  (Other)   :74                                                        
##         trans    drv          cty             hwy        fl     
##  auto(l4)  :83   4:103   Min.   : 9.00   Min.   :12.00   c:  1  
##  manual(m5):58   f:106   1st Qu.:14.00   1st Qu.:18.00   d:  5  
##  auto(l5)  :39   r: 25   Median :17.00   Median :24.00   e:  8  
##  manual(m6):19           Mean   :16.86   Mean   :23.44   p: 52  
##  auto(s6)  :16           3rd Qu.:19.00   3rd Qu.:27.00   r:168  
##  auto(l6)  : 6           Max.   :35.00   Max.   :44.00          
##  (Other)   :13                                                  
##         class   
##  2seater   : 5  
##  compact   :47  
##  midsize   :41  
##  minivan   :11  
##  pickup    :33  
##  subcompact:35  
##  suv       :62
\end{verbatim}

Dengan mengubah tipe variabel menjadi faktor, kategori pada
masing-masing variabel akan terlihat dan dihitung frekuensinya. Tampilan
ini memudahkan Anda untuk mengetahui bahwa produsen dodge dan toyota
memiliki oberservasi terbanyak dibandingkan produsen mobil lainnya.

Bandingkan dengan hasil ringkasan set data aslinya yang kurang
informatif:

\begin{Shaded}
\begin{Highlighting}[]
\CommentTok{# Menampilkan ringkasan umum dari set data df.mobil}
\KeywordTok{summary}\NormalTok{(mpg)}
\end{Highlighting}
\end{Shaded}

\begin{verbatim}
##  manufacturer          model               displ            year     
##  Length:234         Length:234         Min.   :1.600   Min.   :1999  
##  Class :character   Class :character   1st Qu.:2.400   1st Qu.:1999  
##  Mode  :character   Mode  :character   Median :3.300   Median :2004  
##                                        Mean   :3.472   Mean   :2004  
##                                        3rd Qu.:4.600   3rd Qu.:2008  
##                                        Max.   :7.000   Max.   :2008  
##       cyl           trans               drv                 cty       
##  Min.   :4.000   Length:234         Length:234         Min.   : 9.00  
##  1st Qu.:4.000   Class :character   Class :character   1st Qu.:14.00  
##  Median :6.000   Mode  :character   Mode  :character   Median :17.00  
##  Mean   :5.889                                         Mean   :16.86  
##  3rd Qu.:8.000                                         3rd Qu.:19.00  
##  Max.   :8.000                                         Max.   :35.00  
##       hwy             fl               class          
##  Min.   :12.00   Length:234         Length:234        
##  1st Qu.:18.00   Class :character   Class :character  
##  Median :24.00   Mode  :character   Mode  :character  
##  Mean   :23.44                                        
##  3rd Qu.:27.00                                        
##  Max.   :44.00
\end{verbatim}

\textbf{MENGUBAH NAMA VARIABEL}

Nama variabel set data \texttt{mpg} menggunakan singkatan yang pendek
dan berbahasa Inggris. Untuk memudahkan pemahaman, nama variabel
tersebut akan diganti ke bahasa Indonesia serta tanpa singkatan.

\begin{Shaded}
\begin{Highlighting}[]
\CommentTok{# Menampilkan nama variabel set data df.mobil}
\KeywordTok{names}\NormalTok{(df.mobil)}
\end{Highlighting}
\end{Shaded}

\begin{verbatim}
##  [1] "manufacturer" "model"        "displ"        "year"        
##  [5] "cyl"          "trans"        "drv"          "cty"         
##  [9] "hwy"          "fl"           "class"
\end{verbatim}

Perhatikan urutan dari output di atas. Urutan nama variabel dimulai dari
baris paling atas kiri ke kanan, lalu baris kedua mulai kiri ke kanan
dan seterusnya. Angka di dalam kurung kotak \texttt{{[}{]}} adalah alat
bantu untuk mempermudah identifikasi nomor urut variabel. Contohnya
\texttt{{[}7{]}} adalah nomor urut variabel `drv', sebelah kanannya,
`cty' sudah tentu nomor delapan dan seterusnya.

Mengubah nama variabel bisa dilakukan satu per satu atau sekaligus.

Berikut cara untuk mengubah salah satu nama variabel pada sebuah set
data.

\begin{Shaded}
\begin{Highlighting}[]
\CommentTok{# Mengubah nama variabel displ menjadi kapasitas}
\KeywordTok{names}\NormalTok{(df.mobil)[}\DecValTok{3}\NormalTok{] <-}\StringTok{ "kapasitas"}
\end{Highlighting}
\end{Shaded}

Penjelasan dari kode di atas adalah mengubah nama variabel yang
merupakan elemen vektor ketiga set data df.mobil menjadi ``kapasitas''.

Yang perlu Anda identifikasi adalah nomor urut ke berapa variabel yang
akan diubah. Variabel `displ' urutan ketiga. Dengan menggunakan prinsip
elemen vektor, maka digunakan \texttt{{[}3{]}} sebagai akses elemen
dengan urutan yang ketiga.

Pastikan bahwa nama variabel `displ' sudah diubah.

\begin{Shaded}
\begin{Highlighting}[]
\CommentTok{# Menampilkan nama variabel set data df.mobil}
\KeywordTok{names}\NormalTok{(df.mobil)}
\end{Highlighting}
\end{Shaded}

\begin{verbatim}
##  [1] "manufacturer" "model"        "kapasitas"    "year"        
##  [5] "cyl"          "trans"        "drv"          "cty"         
##  [9] "hwy"          "fl"           "class"
\end{verbatim}

Variabel kapasitas sudah ada di set data df.mobil.

Lalu, bagaimana cara mengubah semua nama variabel sekaligus?

\begin{Shaded}
\begin{Highlighting}[]
\CommentTok{# Mengubah seluruh nama variabel }
\KeywordTok{names}\NormalTok{(df.mobil) <-}\StringTok{ }\KeywordTok{c}\NormalTok{(}\StringTok{"produsen"}\NormalTok{, }\StringTok{"model"}\NormalTok{, }\StringTok{"kapasitas"}\NormalTok{, }\StringTok{"tahun"}\NormalTok{, }\StringTok{"silinder"}\NormalTok{, }\StringTok{"transmisi"}\NormalTok{, }\StringTok{"roda.gerak"}\NormalTok{, }\StringTok{"kota"}\NormalTok{, }\StringTok{"jalan.tol"}\NormalTok{, }\StringTok{"jenis.bensin"}\NormalTok{, }\StringTok{"kelas"}\NormalTok{)}
\end{Highlighting}
\end{Shaded}

Dengan menggunakan fungsi \texttt{names()} dan \texttt{c()}, kesebelas
nama varibel baru disusun secara berurutan.

\begin{Shaded}
\begin{Highlighting}[]
\CommentTok{# Menampilkan nama variabel set data df.mobil}
\KeywordTok{names}\NormalTok{(df.mobil)}
\end{Highlighting}
\end{Shaded}

\begin{verbatim}
##  [1] "produsen"     "model"        "kapasitas"    "tahun"       
##  [5] "silinder"     "transmisi"    "roda.gerak"   "kota"        
##  [9] "jalan.tol"    "jenis.bensin" "kelas"
\end{verbatim}

Seluruh nama variabel telah diubah.

\textbf{MENGUBAH ISI DATA}

Konsep mengubah isi data sebenarnya mirip seperti halnya konsep
\emph{``find and replace''} pada aplikasi \emph{spreadsheet}. Untuk
manipulasi ini digunakan fungsi \texttt{gsub()}.

Argumen pertama huruf/kata yang akan diganti, argumen kedua huruf/kata
pengganti, dan argumen terakhir adalah target variabelnya.

\begin{Shaded}
\begin{Highlighting}[]
\CommentTok{# Mengubah isi variabel drv menjadi yang sebelumnya }
\CommentTok{# satu huruf menjadi satu kata berbahasa Indonesia }
\NormalTok{df.mobil}\OperatorTok{$}\NormalTok{roda.gerak <-}\StringTok{ }\KeywordTok{gsub}\NormalTok{(}\StringTok{"f"}\NormalTok{, }\StringTok{"depan"}\NormalTok{, df.mobil}\OperatorTok{$}\NormalTok{roda.gerak)}
\NormalTok{df.mobil}\OperatorTok{$}\NormalTok{roda.gerak <-}\StringTok{ }\KeywordTok{gsub}\NormalTok{(}\StringTok{"r"}\NormalTok{, }\StringTok{"belakang"}\NormalTok{, df.mobil}\OperatorTok{$}\NormalTok{roda.gerak)}
\NormalTok{df.mobil}\OperatorTok{$}\NormalTok{roda.gerak <-}\StringTok{ }\KeywordTok{gsub}\NormalTok{(}\StringTok{"4"}\NormalTok{, }\StringTok{"4roda"}\NormalTok{, df.mobil}\OperatorTok{$}\NormalTok{roda.gerak)}
\end{Highlighting}
\end{Shaded}

Hasil pengubahan isi data ini bisa dilihat menggunakan fungsi
\texttt{table()}.

\begin{Shaded}
\begin{Highlighting}[]
\CommentTok{# Menampilkan ringkasan variabel roda.gerak pada set data df.mobil}
\KeywordTok{table}\NormalTok{(df.mobil}\OperatorTok{$}\NormalTok{roda.gerak)}
\end{Highlighting}
\end{Shaded}

\begin{verbatim}
## 
##    4roda belakang    depan 
##      103       25      106
\end{verbatim}

Fungsi \texttt{table()} menampilkan ringkasan isi berupa banyaknya
frekuensi faktor yang ada di salah satu variabel.

Bandingkan dengan data aslinya.

\begin{Shaded}
\begin{Highlighting}[]
\CommentTok{# Menampilkan ringkasan variabel drv pada set data mpg}
\KeywordTok{table}\NormalTok{(mpg}\OperatorTok{$}\NormalTok{drv)}
\end{Highlighting}
\end{Shaded}

\begin{verbatim}
## 
##   4   f   r 
## 103 106  25
\end{verbatim}

Yang baru saja dilakukan adalah cara mengubah isi data bertipe karakter.
Bagaimana cara mengubah data numerik?

Misalnya, Anda ingin mengubah isi data variabel `kota' dan `jalan.tol'
pada set data df.mobil yang sebelumnya menggunakan satuan miles per
gallon (mpg) menjadi kilometer per liter (untuk seterusnya disingkat
kpl). Satuan kpl dalam perhitungan konsumsi BBM lebih familiar di
Indonesia.

1 miles per gallon (mpg standar Amerika) setara dengan 0.425144
kilometer per liter (kpl). Agar memudahkan perhitungan, dibulatkan
menjadi 0.43 kpl.

Di bab sebelumnya, Anda sudah belajar bagaimana R berfungsi sebagai
\protect\hyperlink{kalkulator-interaktif}{Kalkulator Interaktif}. Ya,
Anda cukup menggunakan operator \texttt{*} untuk mengalikan 0.43
variabel `kota' dan `jalan.tol' sehingga isi data berubah menjadi satuan
kpl.

\begin{Shaded}
\begin{Highlighting}[]
\CommentTok{# Mengubah isi dari variabel kota dan jalan tol menjadi satuan kpl}
\CommentTok{# dengan mengalikan 0.43}
\NormalTok{df.mobil}\OperatorTok{$}\NormalTok{kota <-}\StringTok{ }\NormalTok{df.mobil}\OperatorTok{$}\NormalTok{kota }\OperatorTok{*}\StringTok{ }\FloatTok{0.43}
\NormalTok{df.mobil}\OperatorTok{$}\NormalTok{jalan.tol <-}\StringTok{ }\NormalTok{df.mobil}\OperatorTok{$}\NormalTok{jalan.tol }\OperatorTok{*}\StringTok{ }\FloatTok{0.43}
\end{Highlighting}
\end{Shaded}

Melihat hasilnya menggunakan fungsi \texttt{head()}.

\begin{Shaded}
\begin{Highlighting}[]
\CommentTok{# Melihat 5 data pertama dari variabel kota dan jalan.tol pada set data df.mobil}
\KeywordTok{head}\NormalTok{(df.mobil}\OperatorTok{$}\NormalTok{kota, }\DecValTok{5}\NormalTok{)}
\end{Highlighting}
\end{Shaded}

\begin{verbatim}
## [1] 7.74 9.03 8.60 9.03 6.88
\end{verbatim}

\begin{Shaded}
\begin{Highlighting}[]
\KeywordTok{head}\NormalTok{(df.mobil}\OperatorTok{$}\NormalTok{jalan.tol, }\DecValTok{5}\NormalTok{)}
\end{Highlighting}
\end{Shaded}

\begin{verbatim}
## [1] 12.47 12.47 13.33 12.90 11.18
\end{verbatim}

\subsection{\texorpdfstring{Operator Pipa
(\texttt{\%\textgreater{}\%})}{Operator Pipa (\%\textgreater{}\%)}}\label{operator-pipa}

R adalah bahasa fungsional. Artinya, \emph{syntax} atau fungsi banyak
sekali menggunakan tanda kurung \texttt{(} dan \texttt{)}. Ketika Anda
melakukan manipulasi data yang kompleks, maka konsekuensinya kode juga
banyak mengandung tanda kurung. Fungsi dalam fungsi yang bersarang. Ini
membuat kode Anda sulit dibaca dan dimengerti. Operator pipa
\texttt{\%\textgreater{}\%} bisa menyederhanakannya.

Sebagai contoh, Anda melakukan perhitungan berikut.

\begin{Shaded}
\begin{Highlighting}[]
\KeywordTok{sin}\NormalTok{(}\KeywordTok{exp}\NormalTok{(}\KeywordTok{log}\NormalTok{(}\KeywordTok{sqrt}\NormalTok{(}\DecValTok{9}\NormalTok{))))}
\end{Highlighting}
\end{Shaded}

\begin{verbatim}
## [1] 0.14112
\end{verbatim}

Ada 4 fungsi yang terlibat dan banyak tanda kurung sehingga untuk
mengetiknya pun Anda cukup rumit. Struktur urutan pengetikkan pun
dimulai dari dalam ke luar atau dari kiri ke kanan. Mulai dari
\texttt{sqrt(9)} lalu dikurung oleh \texttt{log()}, dikurung, dikurung
sampai terakhir oleh \texttt{sin()}.

Sekarang bandingkan dengan menggunakan \texttt{\%\textgreater{}\%}.
Hasil perhitungannya sama.

\begin{Shaded}
\begin{Highlighting}[]
\CommentTok{# Inisiasi paket magrittr agar mengenali operator pipa}
\KeywordTok{library}\NormalTok{(magrittr)}

\CommentTok{# Contoh penggunaan operator pipa}
\DecValTok{9} \OperatorTok{%>%}\StringTok{ }\KeywordTok{sqrt}\NormalTok{() }\OperatorTok{%>%}\StringTok{ }\KeywordTok{log}\NormalTok{() }\OperatorTok{%>%}\StringTok{ }\KeywordTok{exp}\NormalTok{() }\OperatorTok{%>%}\StringTok{ }\KeywordTok{sin}\NormalTok{()}
\end{Highlighting}
\end{Shaded}

\begin{verbatim}
## [1] 0.14112
\end{verbatim}

Operator pipa memberikan kemudahan memahami kode pemrograman dalam hal:

\begin{itemize}
\tightlist
\item
  Struktur urutan fungsi dari kiri ke kanan (yang biasanya dari kanan ke
  kiri)
\item
  Menghindari banyaknya tanda kurung yang bersarang dari fungsi
\item
  Menambah fungsi sesuai urutan jika diperlukan
\end{itemize}

Contoh lainnya:

\begin{itemize}
\tightlist
\item
  Kode umum yang biasa digunakan untuk melihat 7 baris pertama pada set
  data df.mobil: \texttt{head(df.mobil,\ 7)}
\item
  Kode menggunakan operator pipa:
  \texttt{df.mobil\ \%\textgreater{}\%\ head(7)}
\end{itemize}

\begin{Shaded}
\begin{Highlighting}[]
\CommentTok{# Contoh lain penggunaan operator pipa}
\CommentTok{# melihat 7 baris pertama pada set data df.mobil}
\NormalTok{df.mobil }\OperatorTok{%>%}\StringTok{ }\KeywordTok{head}\NormalTok{(}\DecValTok{7}\NormalTok{)}
\end{Highlighting}
\end{Shaded}

\begin{verbatim}
## # A tibble: 7 x 11
##   produsen model kapasitas tahun silinder transmisi roda.gerak  kota
##   <fct>    <chr>     <dbl> <fct> <fct>    <fct>     <chr>      <dbl>
## 1 audi     a4          1.8 1999  4        auto(l5)  depan       7.74
## 2 audi     a4          1.8 1999  4        manual(m~ depan       9.03
## 3 audi     a4          2   2008  4        manual(m~ depan       8.6 
## 4 audi     a4          2   2008  4        auto(av)  depan       9.03
## 5 audi     a4          2.8 1999  6        auto(l5)  depan       6.88
## 6 audi     a4          2.8 1999  6        manual(m~ depan       7.74
## 7 audi     a4          3.1 2008  6        auto(av)  depan       7.74
## # ... with 3 more variables: jalan.tol <dbl>, jenis.bensin <fct>,
## #   kelas <fct>
\end{verbatim}

Anda harus mengenal dan memahami operator pipa ini terlebih dulu karena
pembahasan ke depan akan sering digunakan.

\subsection{Memilih Variabel Data}\label{memilih-variabel-data}

Untuk memilih beberapa variabel atau kolom pada sebuah set data, gunakan
fungsi \texttt{select()}.

\begin{Shaded}
\begin{Highlighting}[]
\CommentTok{# Memilih variabel produsen, tahun, dan roda.gerak dari 11 variabel yang ada}
\CommentTok{# pada set data df.mobil}
\NormalTok{contoh.pilih <-}\StringTok{ }\NormalTok{df.mobil }\OperatorTok{%>%}\StringTok{ }
\StringTok{  }\KeywordTok{select}\NormalTok{(produsen, tahun, roda.gerak)}
\end{Highlighting}
\end{Shaded}

\begin{Shaded}
\begin{Highlighting}[]
\CommentTok{# Melihat 8 baris pertama pada set data contoh.pilih}
\NormalTok{contoh.pilih }\OperatorTok{%>%}\StringTok{  }\KeywordTok{head}\NormalTok{(}\DecValTok{8}\NormalTok{)}
\end{Highlighting}
\end{Shaded}

\begin{verbatim}
## # A tibble: 8 x 3
##   produsen tahun roda.gerak
##   <fct>    <fct> <chr>     
## 1 audi     1999  depan     
## 2 audi     1999  depan     
## 3 audi     2008  depan     
## 4 audi     2008  depan     
## 5 audi     1999  depan     
## 6 audi     1999  depan     
## 7 audi     2008  depan     
## 8 audi     1999  4roda
\end{verbatim}

Kode pemrograman di atas dapat digabung tanpa perlu menyimpan hasilnya
di sebuah variabel baru.

\begin{Shaded}
\begin{Highlighting}[]
\CommentTok{# Memilih variabel produsen, tahun, dan roda.gerak }
\CommentTok{# lalu menampilkan 8 baris pertama dari set data df.mobil}
\NormalTok{df.mobil }\OperatorTok{%>%}\StringTok{ }
\StringTok{  }\KeywordTok{select}\NormalTok{(produsen, tahun, roda.gerak) }\OperatorTok{%>%}
\StringTok{  }\KeywordTok{head}\NormalTok{(}\DecValTok{8}\NormalTok{)}
\end{Highlighting}
\end{Shaded}

\begin{verbatim}
## # A tibble: 8 x 3
##   produsen tahun roda.gerak
##   <fct>    <fct> <chr>     
## 1 audi     1999  depan     
## 2 audi     1999  depan     
## 3 audi     2008  depan     
## 4 audi     2008  depan     
## 5 audi     1999  depan     
## 6 audi     1999  depan     
## 7 audi     2008  depan     
## 8 audi     1999  4roda
\end{verbatim}

\subsection{Menyaring Baris Data}\label{menyaring-baris-data}

Fungsi \texttt{filter()} digunakan untuk menyaring baris pada sebuah set
data dengan kondisi tertentu.

Misalnya, Anda ingin mengetahui model mobil yang merupakan kelas dua
tempat duduk atau 2seater dari set data \texttt{df.mobil}.

\begin{Shaded}
\begin{Highlighting}[]
\CommentTok{# Memilih variabel model dan kelas}
\CommentTok{# lalu menyaring baris data dengan kondisi kelas 2seater}
\NormalTok{df.mobil }\OperatorTok{%>%}\StringTok{ }
\StringTok{  }\KeywordTok{select}\NormalTok{(model, kelas) }\OperatorTok{%>%}
\StringTok{  }\KeywordTok{filter}\NormalTok{(kelas }\OperatorTok{==}\StringTok{ "2seater"}\NormalTok{)}
\end{Highlighting}
\end{Shaded}

\begin{verbatim}
## # A tibble: 5 x 2
##   model    kelas  
##   <chr>    <fct>  
## 1 corvette 2seater
## 2 corvette 2seater
## 3 corvette 2seater
## 4 corvette 2seater
## 5 corvette 2seater
\end{verbatim}

\subsection{Mengurutkan Data}\label{mengurutkan-data}

Anda bisa mengurutkan data dengan menggunakan fungsi \texttt{arrange()}.

Anda ingin mengetahui 10 peringkat teratas model mobil yang memiliki
efisiensi konsumsi bensin (kpl) tertinggi dalam kota dari set data
\texttt{df.mobil}.

\begin{Shaded}
\begin{Highlighting}[]
\CommentTok{# Memilih variabel model, transmisi, dan kota}
\CommentTok{# lalu mengurutkan kpl dari yang tertinggi ke terendah}
\CommentTok{# menggunakan fungsi arrange dan desc}
\CommentTok{# serta menampilkan hasilnya untuk 10 baris teratas}
\NormalTok{df.mobil }\OperatorTok{%>%}\StringTok{ }
\StringTok{  }\KeywordTok{select}\NormalTok{(model, transmisi, kota) }\OperatorTok{%>%}
\StringTok{  }\KeywordTok{arrange}\NormalTok{(}\KeywordTok{desc}\NormalTok{(kota)) }\OperatorTok{%>%}\StringTok{ }
\StringTok{  }\KeywordTok{head}\NormalTok{(}\DecValTok{10}\NormalTok{)}
\end{Highlighting}
\end{Shaded}

\begin{verbatim}
## # A tibble: 10 x 3
##    model      transmisi   kota
##    <chr>      <fct>      <dbl>
##  1 new beetle manual(m5)  15.0
##  2 jetta      manual(m5)  14.2
##  3 new beetle auto(l4)    12.5
##  4 civic      manual(m5)  12.0
##  5 corolla    manual(m5)  12.0
##  6 civic      manual(m5)  11.2
##  7 corolla    manual(m5)  11.2
##  8 corolla    auto(l4)    11.2
##  9 civic      manual(m5)  10.8
## 10 civic      auto(l5)    10.8
\end{verbatim}

\subsection{Menambah Variabel Data}\label{menambah-variabel-data}

Pada contoh kali ini, akan ditambahkan variabel atau kolom efisiensi
bensin untuk penggunaan dalam kota dengan satuan kilometer per liter
dari set data \texttt{mpg}. Untuk manipulasi ini, fungsi
\texttt{mutate()} akan digunakan.

\begin{Shaded}
\begin{Highlighting}[]
\CommentTok{# Memilih variabel model, cty, dan hwy}
\CommentTok{# lalu menambahkan variabel baru: kota.kpl dan tol.kpl}
\CommentTok{# yang merupakan konversi satuan mpg ke kpl}
\CommentTok{# serta menampilkan hasilnya untuk 10 baris teratas}
\NormalTok{mpg }\OperatorTok{%>%}\StringTok{ }
\StringTok{  }\KeywordTok{select}\NormalTok{(model, cty, hwy) }\OperatorTok{%>%}
\StringTok{  }\KeywordTok{mutate}\NormalTok{(}\DataTypeTok{kota.kpl =}\NormalTok{ cty }\OperatorTok{*}\StringTok{ }\FloatTok{0.43}\NormalTok{, }\DataTypeTok{tol.kpl =}\NormalTok{ hwy }\OperatorTok{*}\StringTok{ }\FloatTok{0.43}\NormalTok{) }\OperatorTok{%>%}
\StringTok{  }\KeywordTok{head}\NormalTok{(}\DecValTok{10}\NormalTok{)}
\end{Highlighting}
\end{Shaded}

\begin{verbatim}
## # A tibble: 10 x 5
##    model        cty   hwy kota.kpl tol.kpl
##    <chr>      <int> <int>    <dbl>   <dbl>
##  1 a4            18    29     7.74    12.5
##  2 a4            21    29     9.03    12.5
##  3 a4            20    31     8.6     13.3
##  4 a4            21    30     9.03    12.9
##  5 a4            16    26     6.88    11.2
##  6 a4            18    26     7.74    11.2
##  7 a4            18    27     7.74    11.6
##  8 a4 quattro    18    26     7.74    11.2
##  9 a4 quattro    16    25     6.88    10.8
## 10 a4 quattro    20    28     8.6     12.0
\end{verbatim}

\subsection{Mengelompokkan dan Meringkas
Data}\label{mengelompokkan-dan-meringkas-data}

Fungsi manipulasi terakhir yang dibahas adalah fungsi
\texttt{group\_by()} dan \texttt{summarise()} untuk mengelompokkan dan
meringkas data.

Kembali ke set data \texttt{df.mobil}. Anda ingin mengetahui berapa
rata-rata efisiensi konsumsi bensin penggunaan dalam kota dan jalan tol
untuk masing-masing kategori roda gerak.

\begin{Shaded}
\begin{Highlighting}[]
\CommentTok{# Mengelompokkan berdasarkan kategori roda.gerak}
\CommentTok{# lalu menampilkan rata-rata efisiensi menggunakan fungsi mean}
\CommentTok{# untuk masing-masing konsumsi kpl kota dan jalan tol}
\NormalTok{df.mobil }\OperatorTok{%>%}\StringTok{ }
\StringTok{  }\KeywordTok{group_by}\NormalTok{(roda.gerak) }\OperatorTok{%>%}
\StringTok{  }\KeywordTok{summarise}\NormalTok{(}\DataTypeTok{kota.kpl =} \KeywordTok{mean}\NormalTok{(kota), }\DataTypeTok{tol.kpl =} \KeywordTok{mean}\NormalTok{(jalan.tol))}
\end{Highlighting}
\end{Shaded}

\begin{verbatim}
## # A tibble: 3 x 3
##   roda.gerak kota.kpl tol.kpl
##   <chr>         <dbl>   <dbl>
## 1 4roda          6.16    8.25
## 2 belakang       6.05    9.03
## 3 depan          8.59   12.1
\end{verbatim}

\section{Interpretasi Data}\label{interpretasi-data}

Interpretasi data berarti menyimpulkan hasil analisis data yang sudah
Anda lakukan. Anda sudah bersusah payah membuat visualisasi serta
manipulasi data, namun tidak diinterpretasikan, maka hasilnya nol besar.
Anda harus menarik kesimpulan. Anda harus mendapatkan cerita dari data.
Inti dari analisis data yang Anda sudah lakukan adalah untuk mendapatkan
\emph{insight}.

Memang, kemampuan interpretasi setiap orang berbeda dan subjektif dalam
memandang sebuah permasalahan atau dalam menarik informasi dari data.
Untuk mengasah kemampuan interpretasi ini, Anda setidaknya harus
memiliki 3 sikap positif: tekad/semangat \emph{(passion)}, rasa
penasaran \emph{(curiosity)}, dan antusias.

Ketiga sikap positif ini akan sangat membantu Anda, disamping keahlian
teknis dalam analisis data tentunya. Mendapatkan \emph{insight} dari
kumpulan set data merupakan hal yang cukup menantang. Jika keahlian dan
keterampilan teknis sudah semaksimal mungkin dikeluarkan, namun masih
belum juga optimal dalam menggali kesimpulan dari data, sikap semangat
pantang menyerah, rasa penasaran tinggi, dan antusias yang tidak pernah
habis akan menjadi tambahan energi luar biasa bagi Anda.

Selain faktor \emph{soft skill} yang baru saja dibahas, berikut panduan
pertanyaan-pertanyaan teknis sebagai pemicu untuk membantu Anda dalam
proses interpretasi dan menggali informasi serta \emph{insight} dalam
analisis data.

\begin{itemize}
\item
  \textbf{Sebenarnya apa tujuan utama atau motivasi Anda dalam
  menganalisis set data tertentu? Apa inti permasalahan yang Anda atau
  organisasi/perusahaan hadapi?} Ambil contoh set data \texttt{mpg} yang
  merupakan data efisiensi penggunaan bensin pada mobil. Anggaplah dalam
  beberapa bulan ke depan, Anda berencana membeli mobil pertama yang
  merupakan hasil jerih payah kerja Anda selama ini. Anda tidak ingin
  salah dalam membeli mobil pertama impian Anda. \textbf{Kriteria
  pilihan Anda adalah mobil yang paling irit digunakan di dalam kota}.
  Bisa dikatakan, kriteria pemilihan mobil merupakan salah satu motivasi
  Anda untuk menganalisis dan menggali informasi dari set data
  \texttt{mpg}. Anda ingin tahu, tipe atau jenis mobil seperti apa yang
  memiliki efisiensi penggunaan bensin yang tinggi sehingga Anda bisa
  membeli mobil dengan kriteria atau karakterisitik yang sama di
  Indonesia.
\item
  \textbf{Dari beberapa variabel pada suatu set data, kira-kira variabel
  mana saja yang paling berpengaruh terhadap pengambilan keputusan dalam
  menjawab permasalahan Anda?} Berkaitan dengan contoh sebelumnya, Anda
  mencari tahu mobil yang paling irit, maka tentu saja variabel utama
  yang paling berpengaruh adalah variabel `cty' dan `hwy'. Kedua
  variabel tersebut merupakan data miles per gallon penggunaan bensin di
  kota dan jalan tol. Selain kedua variabel tersebut, Anda juga harus
  mencari variabel lain yang berpengaruh terhadap efisiensi penggunaan
  mobil.
\item
  \textbf{Bagaimana cara mendapatkan temuan atau \emph{insight} pada set
  data yang berguna dalam mengatasi permasalahan Anda?} Seringlah
  latihan dan otak-atik data. Untuk mendapatkan \emph{insight},
  pertama-tama Anda harus memahami data yang Anda analisis. Set data
  dapat Anda pahami dengan melakukan banyak visualisasi dan manipulasi
  sehingga Anda bisa menemukan semacam relasi variabel juga pola-pola
  tertentu yang jika terus dilakukan dan diasah akan meningkatkan
  intuisi Anda dalam menarik \emph{insight} dari set data.
\end{itemize}

\textbf{5 TIPE HUBUNGAN ANTAR VARIABEL DALAM SET DATA}

Ketika Anda bermain dengan set data, ada satu titik Anda menemukan
sebuah pola atau keterkaitan antar variabel yang bisa membantu Anda
mendapatkan \emph{insight}. Anda harus mengetahui 5 tipe hubungan antar
variabel dalam set data yang sering ditemukan:

\begin{enumerate}
\def\labelenumi{\arabic{enumi}.}
\tightlist
\item
  Korelasi
\item
  \emph{Trend}
\item
  Distribusi
\item
  \emph{Outlier}
\item
  Perbandingan dan Peringkat
\end{enumerate}

\subsection{Korelasi}\label{korelasi}

Korelasi antar dua variabel bisa diidentifikasi dengan menggunakan
grafik sebar \emph{(scatter plot)}.

Berikut gambar penjelasannya (sumber:
\url{https://statistics.laerd.com/statistical-guides/pearson-correlation-coefficient-statistical-guide.php}):

\includegraphics[width=1\linewidth]{images/pearson-1-small}

\begin{itemize}
\tightlist
\item
  Korelasi positif: peningkatan nilai pada variabel pertama berbanding
  lurus dengan nilai variabel kedua. Variabel A naik, variabel B naik.
\item
  Korelasi negatif: peningkatan nilai pada variabel pertama berbanding
  terbalik dengan nilai variabel kedua. Variabel A naik, variabel B
  turun.
\item
  Tidak ada korelasi: tidak ada kolerasi antara variabel pertama dan
  kedua.
\end{itemize}

Misalnya, Anda akan mengecek apakah ada korelasi antara variabel
`kapasitas' dan `kota' pada set data \texttt{df.mobil}.

\begin{Shaded}
\begin{Highlighting}[]
\CommentTok{# Membuat grafik sebar antara variabel kapasitas dan kota}
\CommentTok{# pada set data df.mobil menggunakan fungsi geom_point()}
\NormalTok{df.mobil }\OperatorTok{%>%}\StringTok{ }\KeywordTok{ggplot}\NormalTok{(}\KeywordTok{aes}\NormalTok{(}\DataTypeTok{x =}\NormalTok{ kota, }\DataTypeTok{y =}\NormalTok{ kapasitas)) }\OperatorTok{+}\StringTok{ }\KeywordTok{geom_point}\NormalTok{()}
\end{Highlighting}
\end{Shaded}

\includegraphics{ebook_files/figure-latex/unnamed-chunk-73-1.pdf}

Anda bisa lihat grafik sebar diatas. Sebarannya berpola korelasi
negatif. Semakin rendah nilai kapasitas mesin, maka efisiensi penggunaan
bensin mobil di dalam kota akan semakin tinggi, demikian sebaliknya.

Selain melalui grafik sebar, koefisien korelasi juga bisa dihitung
dengan menggunakan fungsi \texttt{cor.test()}.

\begin{Shaded}
\begin{Highlighting}[]
\CommentTok{# Mengitung korelasi antara variabel kapasitas dan kota}
\CommentTok{# pada set data df.mobil menggunakan fungsi cor.test()}
  \KeywordTok{cor.test}\NormalTok{(df.mobil}\OperatorTok{$}\NormalTok{kapasitas, df.mobil}\OperatorTok{$}\NormalTok{kota)}
\end{Highlighting}
\end{Shaded}

\begin{verbatim}
## 
##  Pearson's product-moment correlation
## 
## data:  df.mobil$kapasitas and df.mobil$kota
## t = -20.205, df = 232, p-value < 2.2e-16
## alternative hypothesis: true correlation is not equal to 0
## 95 percent confidence interval:
##  -0.8406782 -0.7467508
## sample estimates:
##       cor 
## -0.798524
\end{verbatim}

Nilai koefisien korelasinya \texttt{-0.798524}. Korelasi negatif.

Korelasinya cukup kuat. Kekuatan korelasi antar dua variabel ditentukan
oleh nilai koefisiennya. Semakin mendekati angka \texttt{1} atau
\texttt{-1}, maka korelasinya semakin kuat.

\subsection{\texorpdfstring{\emph{Trend}}{Trend}}\label{trend}

\emph{Trend} dapat didefinisikan sebagai arah atau pergerakan. Bisa naik
atau turun. Biasanya ditunjukkan melalui grafik garis atau grafik
batang.

Berikut contoh grafik trend (sumber:
\url{http://ririnzuliyaningsih.blogspot.com/2015/05/kemiskinan-dan-kesenjangan-pembangunan.html})

\includegraphics[width=1\linewidth]{images/garis-kemiskinan}

Grafik ini menunjukkan \emph{trend} kemiskinan di Indonesia dari tahun
1996 sampai dengan 2012. Trend naik ditunjukkan mulai tahun 1996 - 1999.
Dari tahun 2006 sampai tahun 2012 terlihat penurunan jumlah penduduk
miskin.

Kesimpulannya, terjadi penurunan \emph{trend} kemiskinan di Indonesia
dalam rentang waktu 16 tahun sejak 1996.

Selain \emph{insight} dari \emph{trend}, Anda juga bisa mencermati dari
rentang 16 tahun tersebut, jumlah penduduk miskin tertinggi ada di tahun
berapa?

Tahun 1998.

Lalu, Anda coba gali lagi, ada apa di tahun 1998 yang menyebabkan angka
kemiskinan menjadi paling tinggi dibandingkan tahun lainnya?

Krisis moneter.

Apakah ketika terjadi krisis di Indonesia, selalu diikuti dengan
peningkatan jumlah kemiskinan di tahun yang sama?

Di tahun 2008 terjadi krisis ekonomi, namun tidak diikuti dengan
meningkatnya jumlah penduduk miskin.

Kenapa tidak terjadi peningkatan jumlah kemiskinan di tahun 2008? Apa
bedanya krisis tahun 1998 dengan tahun 2008? Apa yang menyebabkan
krisis? Apakah ada kaitan antara krisis ekonomi dengan kemiskinan?

Dan seterusnya.

Diiringi dengan rasa penasaran serta antusias, munculkan
pertanyaan-pertanyaan yang bisa membuat Anda berpikir kritis. Dengan
cara ini, Anda bisa menggali \emph{insight} dan informasi dari sebuah
set data atau grafik.

\subsection{Distribusi}\label{distribusi}

Distribusi sering digunakan untuk mengetahui nilai terendah, tertinggi,
rata-rata, median, serta rentang dari suatu variabel pada set data.
Grafik histogram biasa dipakai untuk menunjukkan pola distribusi ini.

Sebagai contoh, Anda ingin mengetahui distribusi jumlah mobil
berdasarkan kapasitas mesinnya pada set data \texttt{df.mobil}.

\begin{Shaded}
\begin{Highlighting}[]
\CommentTok{# Membuat grafik histogram kapasitas mesin}
\CommentTok{# pada set data df.mobil menggunakan fungsi geom_histogram()}
\NormalTok{df.mobil }\OperatorTok{%>%}\StringTok{ }\KeywordTok{ggplot}\NormalTok{(}\KeywordTok{aes}\NormalTok{(}\DataTypeTok{x =}\NormalTok{ kapasitas)) }\OperatorTok{+}\StringTok{ }\KeywordTok{geom_histogram}\NormalTok{(}\DataTypeTok{color =} \StringTok{"white"}\NormalTok{)}
\end{Highlighting}
\end{Shaded}

\includegraphics{ebook_files/figure-latex/unnamed-chunk-76-1.pdf}

Anda bisa lihat, jumlah mobil terbanyak pada set data \texttt{df.mobil}
memiliki kapasitas mesin antara 1 - 2.5 liter.

\subsection{\texorpdfstring{\emph{Outlier}}{Outlier}}\label{outlier}

\emph{Outlier} atau pencilan adalah data obervasi yang menyimpang
terlalu jauh atau bernilai ekstrim dibandingkan data lainnya. Data
\emph{outlier} ini bisa terjadi karena beberapa hal, diantaranya
kesalahan dalam pemasukan data, kesalahan dalam pengambilan sampel, atau
memang ada data -- data ekstrim yang tidak dapat dihindarkan
keberadaannya. Untuk mengidentifikasinya dapat dilakukan dengan beberapa
cara antara lain melalui boxplot atau scatterplot.

Ilustrasi \emph{outlier} (sumber:
\url{http://disabilitydunktank.com/tag/social-security-disability-fraud/}):

\includegraphics[width=1\linewidth]{images/outlier}

Misalkan Anda seorang penyidik yang bekerja di KPK atau kepolisian. Anda
bisa mendapatkan data pendapatan yang diterima pejabat-pejabat Pegawai
Negeri Sipil (PNS) tertentu dari bank. Setelah Anda analisis data
tersebut, Anda menemukan bahwa rata-rata pendapatan pejabat PNS adalah
sekitar 100 juta rupiah per bulan. Anda gali lebih dalam lagi, ternyata
ada 3 pejabat yang memiliki pendapatan 10 miliar rupiah per bulan.
Inilah yang dinamakan \emph{outlier}.

Apa yang harus dilakukan dengan data \emph{outlier} ini? Tergantung
tujuan Anda. Jika Anda seorang penyidik seperti contoh di atas, tentu
Anda memfokuskan pada 3 pejabat tersebut. Data \emph{outlier} ini dapat
dijadikan sebuah \emph{insight} bernilai bagi Anda. Anda bisa melakukan
penyidikan lebih lanjut terhadap ketiga pejabat tersebut. Dari mana
pendapatan mereka berasal?

Jika Anda, misalkan seorang penentu kebijakan di pemerintahan dan ingin
mengkaji usulan kenaikan gaji PNS, maka data \emph{outlier} tersebut
sebaiknya disingkirkan. \emph{Outlier} dapat menyebabkan penyimpangan
hasil analisis data yang Anda lakukan.

\subsection{Perbandingan dan
Peringkat}\label{perbandingan-dan-peringkat}

Cara ini sering dilakukan jika Anda ingin mengetahui nilai tertinggi dan
terendah serta perbandingan antara satu data obervasi dengan yang
lainnya. Yang sering digunakan biasanya grafik batang.

Perbandingan dan Peringkat ini sudah disinggung sebelumnya di bagian
(\#mengurutkan-data).

Langsung ke contoh. Misalkan Anda ingin mengetahui peringkat produsen
mobil yang memiliki efisiensi penggunaan bensin (kpl) tertinggi di jalan
tol dari set data \texttt{df.mobil}.

\begin{Shaded}
\begin{Highlighting}[]
\CommentTok{# Mengelompokkan berdasarkan variabel produsen}
\CommentTok{# lalu meringkas nilai rata-rata efisiensi di jalan tol}
\CommentTok{# dan menampilkannya dalam grafik batang}
\NormalTok{df.mobil }\OperatorTok{%>%}\StringTok{ }
\StringTok{  }\KeywordTok{group_by}\NormalTok{(produsen) }\OperatorTok{%>%}
\StringTok{  }\KeywordTok{summarise}\NormalTok{(}\DataTypeTok{rerata.tol =} \KeywordTok{mean}\NormalTok{(jalan.tol)) }\OperatorTok{%>%}
\StringTok{  }\KeywordTok{ggplot}\NormalTok{(}\KeywordTok{aes}\NormalTok{(}\DataTypeTok{y =}\NormalTok{ rerata.tol, }\DataTypeTok{x =} \KeywordTok{reorder}\NormalTok{(produsen, rerata.tol))) }\OperatorTok{+}
\StringTok{  }\KeywordTok{geom_bar}\NormalTok{(}\DataTypeTok{stat =} \StringTok{"identity"}\NormalTok{, }\DataTypeTok{fill =} \StringTok{"#8B0000"}\NormalTok{) }\OperatorTok{+}\StringTok{ }
\StringTok{  }\KeywordTok{labs}\NormalTok{(}\DataTypeTok{x =} \StringTok{"Produsen"}\NormalTok{, }\DataTypeTok{y =} \StringTok{"Rata-rata Efisiensi (kpl)"}\NormalTok{) }\OperatorTok{+}
\StringTok{  }\KeywordTok{ggtitle}\NormalTok{(}\StringTok{"Peringkat Efisiensi Penggunaan Bensin di Jalan Tol"}\NormalTok{, }
          \DataTypeTok{subtitle =} \StringTok{"Berdasarkan Produsen Mobil"}\NormalTok{) }\OperatorTok{+}
\StringTok{  }\KeywordTok{coord_flip}\NormalTok{()}
\end{Highlighting}
\end{Shaded}

\includegraphics{ebook_files/figure-latex/unnamed-chunk-78-1.pdf}

Dengan melihat grafik batang di atas, Anda dapat mengetahui produsen
yang memiliki efisiensi tertinggi adalah Honda, sedangkan yang terendah
adalah Land Rover. Selain itu, Anda juga bisa membandingkan efisiensi
dari masing-masing produsen mobil lainnya.

\chapter{Latihan, Latihan, Latihan}\label{latihan-latihan-latihan}

Ya, tidak ada yang spesial atau hal yang \emph{wah} di langkah ketiga
ini. Untuk bisa mahir analisis data, Anda harus banyak latihan.

Baca kembali 2 langkah sebelumnya, terutama langkah kedua. Kuasai teknik
dasar analisis terlebih dahulu.

Jalankan kode programnya, \emph{copy paste} boleh, pelajari cara kerja
dan konsepnya. Setelah Anda sudah tidak merasa asing dengan sistem
pemrograman R, coba Anda ketik kode programnya, tidak \emph{copy paste}.
Dengan mengetik, Anda akan lebih cepat terbiasa dengan beberapa fungsi
yang sering digunakan untuk mengolah data.

Lalu, Anda coba latihan dengan set data yang lain. Jelajahi set data
tersebut dengan fungsi-fungsi yang sudah dibahas di buku ini.

Berikut akan disajikan bagaimana cara melakukan analisis data dari mulai
proses pengambilan sampai dengan mendapatkan \emph{insight}.
Fungsi-fungsi yang digunakan pada contoh kali ini sebagian sudah dibahas
di bagian sebelumnya dan sebagian akan dibahas beriringan dengan
penulisan kode.

\section{Unduh, Inisiasi, dan Mempersiapkan Set
Data}\label{unduh-inisiasi-dan-mempersiapkan-set-data}

Set data yang akan digunakan dapat diunduh secara manual pada tautan
berikut:

\begin{itemize}
\tightlist
\item
  \href{http://data.jakarta.go.id/dataset/jumlahpendudukberdasarkanjeniskelamindankewarganegaraandkijakarta/resource/5a948e1c-facd-4eec-8943-bef5c45d16d8}{Data
  Jumlah Penduduk Jakarta Berdasarkan Jenis Kelamin dan Kewarganegaraan}
\end{itemize}

Anda bisa langsung inisiasi dan menyimpan set data ini ke dalam variabel
tanpa harus mengunduhnya terlebih dahulu.

\begin{Shaded}
\begin{Highlighting}[]
\CommentTok{# Meyimpan alamat tautan (URL) data penduduk Jakarta berdasarkan jenis kelamin dan kewarganegaraan dalam variabel alamat}
\NormalTok{alamat <-}\StringTok{ "http://data.jakarta.go.id/dataset/aed0be8f-b36e-457c-9210-6f5e970b987c/resource/5a948e1c-facd-4eec-8943-bef5c45d16d8/download/Data-Jumlah-Penduduk-Berdasarkan-Kewarganegaraan-dan-Jenis-Kelamin-Tahun-2017.csv"}

\CommentTok{# Membaca file csv dari set data dan memasukkannya ke dalam variabel}
\NormalTok{df.penduduk <-}\StringTok{ }\KeywordTok{read.csv}\NormalTok{(alamat, }\DataTypeTok{sep =} \StringTok{","}\NormalTok{)}
\end{Highlighting}
\end{Shaded}

Set data kependudukan Jakarta tersedia dalam bentuk file csv, sehingga
untuk membaca data tersebut digunakan fungsi \texttt{read.csv}.

Melihat struktur dari set data \texttt{df.penduduk}.

\begin{Shaded}
\begin{Highlighting}[]
\KeywordTok{str}\NormalTok{(df.penduduk)}
\end{Highlighting}
\end{Shaded}

\begin{verbatim}
## 'data.frame':    1068 obs. of  8 variables:
##  $ tahun          : int  2017 2017 2017 2017 2017 2017 2017 2017 2017 2017 ...
##  $ provinsi       : Factor w/ 1 level "DKI Jakarta": 1 1 1 1 1 1 1 1 1 1 ...
##  $ kabupaten      : Factor w/ 6 levels "JAKARTA BARAT",..: 6 6 6 6 6 6 2 2 2 2 ...
##  $ kecamatan      : Factor w/ 44 levels "CAKUNG","CEMPAKA PUTIH",..: 22 22 22 21 21 21 9 9 9 9 ...
##  $ kelurahan      : Factor w/ 267 levels "ANCOL","ANGKE",..: 165 164 163 168 167 166 55 24 195 194 ...
##  $ kewarganegaraan: Factor w/ 2 levels "WNA","WNI": 2 2 2 2 2 2 2 2 2 2 ...
##  $ jenis_kelamin  : Factor w/ 2 levels "Laki-laki","Perempuan": 1 1 1 1 1 1 1 1 1 1 ...
##  $ jumlah         : int  3426 3546 1241 1204 2797 1688 1645 9802 11224 9333 ...
\end{verbatim}

Seperti yang Anda lihat, beberapa variabel seperti provinsi, kabupaten,
kecamatan, kelurahan, kewarganegaraan, serta jenis kelamin sudah bertipe
faktor, sehingga tidak perlu diubah lagi. Variabel yang lain seperti
tahun dan jumlah sudah benar bertipe integer.

Yang perlu sedikit diubah adalah penamaan variabel kabupaten. Di DKI
Jakarta, penamaan wilayah administratif setelah propinsi lebih tepatnya
kota, bukan kabupaten.

Masih ingat bagaimana cara mengubah nama variabel?

\begin{Shaded}
\begin{Highlighting}[]
\CommentTok{# Mengubah nama variabel kabupaten menjadi kota pada set data df.penduduk}
\KeywordTok{names}\NormalTok{(df.penduduk)[}\DecValTok{3}\NormalTok{] <-}\StringTok{ "kota"}
\end{Highlighting}
\end{Shaded}

Bagaimana dengan informasi jumlah penduduk Jakarta?

\begin{Shaded}
\begin{Highlighting}[]
\CommentTok{# Jumlah total penduduk Jakarta}
\KeywordTok{sum}\NormalTok{(df.penduduk}\OperatorTok{$}\NormalTok{jumlah)}
\end{Highlighting}
\end{Shaded}

\begin{verbatim}
## [1] 10348570
\end{verbatim}

\begin{Shaded}
\begin{Highlighting}[]
\CommentTok{# Jumlah penduduk Jakarta terbanyak berdasarkan kecamatan}
\NormalTok{df.penduduk }\OperatorTok{%>%}
\StringTok{  }\KeywordTok{select}\NormalTok{(kecamatan, jumlah) }\OperatorTok{%>%}
\StringTok{  }\KeywordTok{group_by}\NormalTok{(kecamatan) }\OperatorTok{%>%}
\StringTok{  }\KeywordTok{summarise}\NormalTok{(}\DataTypeTok{jumlah =} \KeywordTok{sum}\NormalTok{(jumlah)) }\OperatorTok{%>%}
\StringTok{  }\KeywordTok{arrange}\NormalTok{(}\OperatorTok{-}\NormalTok{jumlah) }\OperatorTok{%>%}
\StringTok{  }\KeywordTok{head}\NormalTok{(}\DecValTok{5}\NormalTok{)}
\end{Highlighting}
\end{Shaded}

\begin{verbatim}
## # A tibble: 5 x 2
##   kecamatan     jumlah
##   <fct>          <int>
## 1 CENGKARENG    513064
## 2 CAKUNG        509194
## 3 DUREN SAWIT   406998
## 4 CILINCING     403028
## 5 TANJUNG PRIOK 395022
\end{verbatim}

\begin{Shaded}
\begin{Highlighting}[]
\CommentTok{# Jumlah penduduk Jakarta paling sedikit berdasarkan kecamatan}
\NormalTok{df.penduduk }\OperatorTok{%>%}
\StringTok{  }\KeywordTok{select}\NormalTok{(kecamatan, jumlah) }\OperatorTok{%>%}
\StringTok{  }\KeywordTok{group_by}\NormalTok{(kecamatan) }\OperatorTok{%>%}
\StringTok{  }\KeywordTok{summarise}\NormalTok{(}\DataTypeTok{jumlah =} \KeywordTok{sum}\NormalTok{(jumlah)) }\OperatorTok{%>%}
\StringTok{  }\KeywordTok{arrange}\NormalTok{(jumlah) }\OperatorTok{%>%}
\StringTok{  }\KeywordTok{head}\NormalTok{(}\DecValTok{5}\NormalTok{)}
\end{Highlighting}
\end{Shaded}

\begin{verbatim}
## # A tibble: 5 x 2
##   kecamatan       jumlah
##   <fct>            <int>
## 1 KEP. SERIBU SLT  11342
## 2 KEP. SERIBU UTR  16249
## 3 MENTENG          92026
## 4 CEMPAKA PUTIH    99263
## 5 GAMBIR          101832
\end{verbatim}

\begin{Shaded}
\begin{Highlighting}[]
\CommentTok{# Peringkat jumlah penduduk Jakarta berdasarkan kota}
\NormalTok{df.penduduk }\OperatorTok{%>%}
\StringTok{  }\KeywordTok{select}\NormalTok{(kota, jumlah) }\OperatorTok{%>%}
\StringTok{  }\KeywordTok{group_by}\NormalTok{(kota) }\OperatorTok{%>%}
\StringTok{  }\KeywordTok{summarise}\NormalTok{(}\DataTypeTok{jumlah =} \KeywordTok{sum}\NormalTok{(jumlah)) }\OperatorTok{%>%}
\StringTok{  }\KeywordTok{arrange}\NormalTok{(}\OperatorTok{-}\NormalTok{jumlah)}
\end{Highlighting}
\end{Shaded}

\begin{verbatim}
## # A tibble: 6 x 2
##   kota                jumlah
##   <fct>                <int>
## 1 JAKARTA TIMUR      2946926
## 2 JAKARTA BARAT      2327258
## 3 JAKARTA SELATAN    2190919
## 4 JAKARTA UTARA      1716591
## 5 JAKARTA PUSAT      1139285
## 6 KAB.ADM.KEP.SERIBU   27591
\end{verbatim}

\section{Visualisasi Grafik}\label{visualisasi-grafik}

Set data kependudukan \texttt{df.penduduk} ini dari awal sudah lumayan
rapih dan tidak diperlukan proses manipulasi. Langkah berikutnya
menggunakan \texttt{ggplot} untuk memvisualisasikan data agar bisa
mendapatkan \texttt{insight} dari set data kependudukan ini.

\begin{Shaded}
\begin{Highlighting}[]
\CommentTok{# Inisiasi paket ggplot dan dplyr}
\KeywordTok{library}\NormalTok{(ggplot2)}
\KeywordTok{library}\NormalTok{(dplyr)}

\NormalTok{df.penduduk }\OperatorTok{%>%}\StringTok{ }
\StringTok{  }\KeywordTok{ggplot}\NormalTok{(}\KeywordTok{aes}\NormalTok{(}\DataTypeTok{x =} \KeywordTok{reorder}\NormalTok{(kecamatan, jumlah, sum), }\DataTypeTok{y =}\NormalTok{ jumlah)) }\OperatorTok{+}
\StringTok{  }\KeywordTok{geom_bar}\NormalTok{(}\DataTypeTok{stat =} \StringTok{"identity"}\NormalTok{, }\DataTypeTok{fill =} \StringTok{"#8B0000"}\NormalTok{) }\OperatorTok{+}
\StringTok{  }\KeywordTok{coord_flip}\NormalTok{() }\OperatorTok{+}
\StringTok{  }\KeywordTok{theme}\NormalTok{(}\DataTypeTok{axis.text.y =} \KeywordTok{element_text}\NormalTok{(}\DataTypeTok{size =} \DecValTok{6}\NormalTok{)) }\OperatorTok{+}
\StringTok{  }\KeywordTok{scale_y_continuous}\NormalTok{(}\DataTypeTok{labels =}\NormalTok{ scales}\OperatorTok{::}\KeywordTok{comma_format}\NormalTok{()) }\OperatorTok{+}
\StringTok{  }\KeywordTok{labs}\NormalTok{(}\DataTypeTok{x =} \StringTok{"Kecamatan"}\NormalTok{, }\DataTypeTok{y =} \StringTok{"Jumlah Penduduk"}\NormalTok{) }\OperatorTok{+}
\StringTok{  }\KeywordTok{ggtitle}\NormalTok{(}\StringTok{"Perbandingan Jumlah Penduduk Jakarta Berdasarkan Kecamatan"}\NormalTok{, }
          \DataTypeTok{subtitle =} \StringTok{"Tahun 2017"}\NormalTok{)}
\end{Highlighting}
\end{Shaded}

\includegraphics{ebook_files/figure-latex/unnamed-chunk-86-1.pdf}

Penjelasan kode:

\begin{itemize}
\item
  \texttt{ggplot(aes(x\ =\ reorder(kecamatan,\ jumlah,\ sum),\ y\ =\ jumlah))}:
  mengurutkan faktor kecamatan berdasarkan jumlah dengan fungsi
  \texttt{reorder}.
\item
  \texttt{theme(axis.text.y\ =\ element\_text(size\ =\ 6))}: mengatur
  huruf pada axis y menjadi berukuran 6.
\item
  \texttt{scale\_y\_continuous(labels\ =\ scales::comma\_format()}:
  mengatur skala angka menjadi format yang mudah dibaca.
\end{itemize}

\begin{Shaded}
\begin{Highlighting}[]
\NormalTok{df.penduduk }\OperatorTok{%>%}\StringTok{ }
\StringTok{  }\KeywordTok{ggplot}\NormalTok{(}\KeywordTok{aes}\NormalTok{(}\DataTypeTok{x =} \KeywordTok{reorder}\NormalTok{(kota, jumlah, sum), }\DataTypeTok{y =}\NormalTok{ jumlah)) }\OperatorTok{+}
\StringTok{  }\KeywordTok{geom_bar}\NormalTok{(}\DataTypeTok{stat =} \StringTok{"identity"}\NormalTok{, }\DataTypeTok{fill =} \StringTok{"#8B0000"}\NormalTok{) }\OperatorTok{+}
\StringTok{  }\KeywordTok{coord_flip}\NormalTok{() }\OperatorTok{+}
\StringTok{  }\KeywordTok{theme}\NormalTok{(}\DataTypeTok{axis.text.y =} \KeywordTok{element_text}\NormalTok{(}\DataTypeTok{size =} \DecValTok{8}\NormalTok{)) }\OperatorTok{+}
\StringTok{  }\KeywordTok{theme}\NormalTok{(}\DataTypeTok{axis.text.x =} \KeywordTok{element_text}\NormalTok{(}\DataTypeTok{size =} \DecValTok{7}\NormalTok{)) }\OperatorTok{+}
\StringTok{  }\KeywordTok{scale_y_continuous}\NormalTok{(}\DataTypeTok{labels =}\NormalTok{ scales}\OperatorTok{::}\KeywordTok{comma_format}\NormalTok{()) }\OperatorTok{+}
\StringTok{  }\KeywordTok{labs}\NormalTok{(}\DataTypeTok{x =} \StringTok{"Kota"}\NormalTok{, }\DataTypeTok{y =} \StringTok{"Jumlah Penduduk"}\NormalTok{) }\OperatorTok{+}
\StringTok{  }\KeywordTok{ggtitle}\NormalTok{(}\StringTok{"Perbandingan Jumlah Penduduk Jakarta Berdasarkan Kota"}\NormalTok{,}
          \DataTypeTok{subtitle =} \StringTok{"Tahun 2017"}\NormalTok{)}
\end{Highlighting}
\end{Shaded}

\includegraphics{ebook_files/figure-latex/unnamed-chunk-87-1.pdf}

\begin{Shaded}
\begin{Highlighting}[]
\NormalTok{df.penduduk }\OperatorTok{%>%}\StringTok{ }
\StringTok{  }\KeywordTok{ggplot}\NormalTok{(}\KeywordTok{aes}\NormalTok{(}\DataTypeTok{x =} \KeywordTok{reorder}\NormalTok{(kecamatan, jumlah, sum), }\DataTypeTok{y =}\NormalTok{ jumlah)) }\OperatorTok{+}
\StringTok{  }\KeywordTok{geom_bar}\NormalTok{(}\DataTypeTok{stat =} \StringTok{"identity"}\NormalTok{, }\KeywordTok{aes}\NormalTok{(}\DataTypeTok{fill =}\NormalTok{ jenis_kelamin)) }\OperatorTok{+}
\StringTok{  }\KeywordTok{coord_flip}\NormalTok{() }\OperatorTok{+}
\StringTok{  }\KeywordTok{theme}\NormalTok{(}\DataTypeTok{axis.text.y =} \KeywordTok{element_text}\NormalTok{(}\DataTypeTok{size =} \DecValTok{6}\NormalTok{)) }\OperatorTok{+}
\StringTok{  }\KeywordTok{scale_y_continuous}\NormalTok{(}\DataTypeTok{labels =}\NormalTok{ scales}\OperatorTok{::}\KeywordTok{comma_format}\NormalTok{()) }\OperatorTok{+}
\StringTok{  }\KeywordTok{labs}\NormalTok{(}\DataTypeTok{x =} \StringTok{"Kecamatan"}\NormalTok{, }\DataTypeTok{y =} \StringTok{"Jumlah Penduduk"}\NormalTok{) }\OperatorTok{+}
\StringTok{  }\KeywordTok{ggtitle}\NormalTok{(}\StringTok{"Perbandingan Jumlah Penduduk Jakarta"}\NormalTok{,}
          \DataTypeTok{subtitle =} \StringTok{"Berdasarkan Jenis Kelamin Per Kecamatan Tahun 2017"}\NormalTok{)}
\end{Highlighting}
\end{Shaded}

\includegraphics{ebook_files/figure-latex/unnamed-chunk-88-1.pdf}

Argumen \texttt{fill} pada kode di atas berfungsi untuk menampilkan
atribut estetik berdasarkan variabel jenis\_kelamin, sehingga pada
grafik terlihat komposisi penduduk laki-laki dan perempuan yang berbeda
warna.

\begin{Shaded}
\begin{Highlighting}[]
\NormalTok{df.penduduk }\OperatorTok{%>%}\StringTok{ }
\StringTok{  }\KeywordTok{ggplot}\NormalTok{(}\KeywordTok{aes}\NormalTok{(}\DataTypeTok{x =} \KeywordTok{reorder}\NormalTok{(kecamatan, jumlah, sum), }\DataTypeTok{y =}\NormalTok{ jumlah)) }\OperatorTok{+}
\StringTok{  }\KeywordTok{geom_bar}\NormalTok{(}\DataTypeTok{stat =} \StringTok{"identity"}\NormalTok{, }\DataTypeTok{fill =} \StringTok{"#8B0000"}\NormalTok{) }\OperatorTok{+}
\StringTok{  }\KeywordTok{coord_flip}\NormalTok{() }\OperatorTok{+}
\StringTok{  }\KeywordTok{theme}\NormalTok{(}\DataTypeTok{axis.text.y =} \KeywordTok{element_text}\NormalTok{(}\DataTypeTok{size =} \DecValTok{6}\NormalTok{)) }\OperatorTok{+}
\StringTok{  }\KeywordTok{scale_y_continuous}\NormalTok{(}\DataTypeTok{labels =}\NormalTok{ scales}\OperatorTok{::}\KeywordTok{comma_format}\NormalTok{()) }\OperatorTok{+}
\StringTok{  }\KeywordTok{facet_wrap}\NormalTok{(}\OperatorTok{~}\NormalTok{jenis_kelamin) }\OperatorTok{+}
\StringTok{  }\KeywordTok{labs}\NormalTok{(}\DataTypeTok{x =} \StringTok{"Kecamatan"}\NormalTok{, }\DataTypeTok{y =} \StringTok{"Jumlah Penduduk"}\NormalTok{) }\OperatorTok{+}
\StringTok{  }\KeywordTok{ggtitle}\NormalTok{(}\StringTok{"Perbandingan Jumlah Penduduk Jakarta"}\NormalTok{,}
          \DataTypeTok{subtitle =} \StringTok{"Berdasarkan Jenis Kelamin Per Kecamatan Tahun 2017"}\NormalTok{)}
\end{Highlighting}
\end{Shaded}

\includegraphics{ebook_files/figure-latex/unnamed-chunk-89-1.pdf}

Fungsi \texttt{facet\_wrap()} membuat grafik terbagi menjadi dua kolom
berdasarkan jenis kelamin. Opsi ini mempermudah pengguna untuk melihat
lebih detail untuk masing-masing kategori.

\section{Visualisasi Peta}\label{visualisasi-peta}

Visualisasi data juga dapat ditampilkan dalam bentuk peta. Selain cukup
informatif, visualisasi peta juga secara estetis indah dipandang dan
keren.

Untuk membuat visualisasi peta Jakarta, Anda perlu menginstal paket
berikut terlebih dahulu.

\begin{Shaded}
\begin{Highlighting}[]
\KeywordTok{install.packages}\NormalTok{(}\StringTok{"devtools"}\NormalTok{)}
\KeywordTok{install.packages}\NormalTok{(}\StringTok{"sf"}\NormalTok{)}
\NormalTok{devtools}\OperatorTok{::}\KeywordTok{install_github}\NormalTok{(}\StringTok{"tidyverse/ggplot2"}\NormalTok{)}
\NormalTok{devtools}\OperatorTok{::}\KeywordTok{install_github}\NormalTok{(}\StringTok{"rasyidstat/indonesia"}\NormalTok{) }
\end{Highlighting}
\end{Shaded}

Inisiasi paket dan persiapan visualiasi peta Jakarta.

\begin{Shaded}
\begin{Highlighting}[]
\CommentTok{# Persiapan data spasial peta Jakarta dengan pembagian wilayah kelurahan}
\NormalTok{kelurahan.jkt <-}\StringTok{ }\KeywordTok{id_map}\NormalTok{(}\StringTok{"jakarta"}\NormalTok{, }\StringTok{"kelurahan"}\NormalTok{)}

\CommentTok{# Menampilkan struktur set data kelurahan.jkt}
\KeywordTok{str}\NormalTok{(kelurahan.jkt)}
\end{Highlighting}
\end{Shaded}

\begin{verbatim}
## Classes 'sf' and 'data.frame':   377 obs. of  9 variables:
##  $ kode_provinsi : int  31 31 31 31 31 31 31 31 31 31 ...
##  $ nama_provinsi : Factor w/ 1 level "DAERAH KHUSUS IBUKOTA JAKARTA": 1 1 1 1 1 1 1 1 1 1 ...
##  $ kode_kota     : int  3171 3172 3172 3172 3172 3171 3172 3172 3171 3172 ...
##  $ nama_kota     : Factor w/ 6 levels "Jakarta Barat",..: 3 4 4 4 4 3 4 4 3 4 ...
##  $ kode_kecamatan: int  3171010 3172020 3172030 3172030 3172010 3171010 3172010 3172020 3171010 3172030 ...
##  $ nama_kecamatan: Factor w/ 45 levels "Cakung","Cempaka Putih",..: 11 7 6 6 34 11 34 7 11 6 ...
##  $ kode_kelurahan: num  3.17e+09 3.17e+09 3.17e+09 3.17e+09 3.17e+09 ...
##  $ nama_kelurahan: Factor w/ 267 levels "Ancol","Angke",..: 35 23 202 162 179 236 87 122 25 34 ...
##  $ geometry      :sfc_POLYGON of length 377; first list element: List of 1
##   ..$ : num [1:1283, 1:2] 107 107 107 107 107 ...
##   ..- attr(*, "class")= chr  "XY" "POLYGON" "sfg"
##  - attr(*, "sf_column")= chr "geometry"
##  - attr(*, "agr")= Factor w/ 3 levels "constant","aggregate",..: NA NA NA NA NA NA NA NA
##   ..- attr(*, "names")= chr  "kode_provinsi" "nama_provinsi" "kode_kota" "nama_kota" ...
\end{verbatim}

Melalui fungsi \texttt{str()} diketahui bahwa set data
\texttt{kelurahan.jkt} bertipe `sf' dan `data.frame'. Simple features
atau sf adalah bentuk set data yang didalamnya berisi variabel geometri
dua dimensi (titik, garis, poligon). Di dalam set data
\texttt{kelurahan.jkt} diwakili oleh variabel `geometry'.

Untuk memvisualisasikan peta, gunakan paket \texttt{ggplot} diikuti
dengan fungsi \texttt{geom\_sf()}.

\begin{Shaded}
\begin{Highlighting}[]
\CommentTok{# Visualisasi peta dasar Jakarta}
\NormalTok{kelurahan.jkt }\OperatorTok{%>%}\StringTok{ }\KeywordTok{ggplot}\NormalTok{() }\OperatorTok{+}\StringTok{ }\KeywordTok{geom_sf}\NormalTok{()}
\end{Highlighting}
\end{Shaded}

\includegraphics{ebook_files/figure-latex/unnamed-chunk-93-1.pdf}

Visualisasi peta dasar Jakarta yang dihasilkan dari kode di atas tidak
terlalu bagus. Seperti yang Anda lihat, Kepulauan Seribu termasuk dalam
wilayah Jakarta, sehingga visualisasi yang dihasilkan mencakup seluruh
wilayah administratif sesuai variabel pada set data
\texttt{kelurahan.jkt}. Agar visualisasi yang dihasilkan lebih estetis,
maka wilayah Kepulauan Seribu akan dikeluarkan dari set data.

\begin{Shaded}
\begin{Highlighting}[]
\CommentTok{# Mengeluarkan data observasi kecamatan KEPULAUAN SERIBU}
\CommentTok{# dan memvisualisasikan peta dengan ggplot dan fungsi geom_sf}
\NormalTok{kelurahan.jkt }\OperatorTok{%>%}
\StringTok{  }\KeywordTok{filter}\NormalTok{(}\OperatorTok{!}\NormalTok{nama_kecamatan }\OperatorTok{%in%}\StringTok{ }\KeywordTok{c}\NormalTok{(}\StringTok{"KEPULAUAN SERIBU UTARA"}\NormalTok{, }\StringTok{"KEPULAUAN SERIBU SELATAN"}\NormalTok{)) }\OperatorTok{%>%}
\StringTok{  }\KeywordTok{ggplot}\NormalTok{() }\OperatorTok{+}\StringTok{ }\KeywordTok{geom_sf}\NormalTok{()}
\end{Highlighting}
\end{Shaded}

\includegraphics{ebook_files/figure-latex/unnamed-chunk-94-1.pdf}

Nah, hasilnya sudah lumayan bagus. Peta difokuskan pada wilayah Jakarta
yang ada di pulau Jawa, tidak termasuk wilayah Kepulauan Seribu.

Selanjutnya yang harus dilakukan adalah menggabungkan set data
\texttt{kelurahan.jkt} dan \texttt{df.penduduk} sehingga visualisasi
peta kepadatan penduduk Jakarta dapat ditampilkan.

Syarat agar dua set data bisa digabung: ada satu variabel memiliki isi
yang sama pada kedua set data. Pada contoh ini, variabel `jumlah' pada
set data \texttt{df.penduduk} akan digabungkan ke set data
\texttt{kelurahan.jkt} melalui variabel kelurahan/nama\_kelurahan.

Jadi, harus dipastikan dulu penamaan variabel kelurahan/nama\_kelurahan
beserta isinya pada kedua set data harus sama.

\begin{Shaded}
\begin{Highlighting}[]
\CommentTok{# Meringkas variabel kelurahan dan jumlah penduduk pada set data df.penduduk}
\CommentTok{# lalu menyimpannya di variabel jml.penduduk}
\NormalTok{jml.penduduk <-}\StringTok{ }
\StringTok{  }\NormalTok{df.penduduk }\OperatorTok{%>%}\StringTok{ }
\StringTok{    }\KeywordTok{group_by}\NormalTok{(kelurahan) }\OperatorTok{%>%}
\StringTok{    }\KeywordTok{summarise}\NormalTok{(}\DataTypeTok{jumlah =} \KeywordTok{sum}\NormalTok{(jumlah))}
\end{Highlighting}
\end{Shaded}

\begin{Shaded}
\begin{Highlighting}[]
\CommentTok{# Mengambil isi variabel nama_kelurahan pada set data keluahan.jkt}
\CommentTok{# dan menyimpannya di variabel kelurahan}
\NormalTok{kelurahan <-}\StringTok{ }\KeywordTok{sort}\NormalTok{(}\KeywordTok{unique}\NormalTok{(kelurahan.jkt}\OperatorTok{$}\NormalTok{nama_kelurahan))}
\end{Highlighting}
\end{Shaded}

\begin{Shaded}
\begin{Highlighting}[]
\CommentTok{# Menggabungkan set data jml.penduduk dan vektor kelurahan}
\CommentTok{# menggunakan fungsi data.frame}
\NormalTok{cek.kelurahan <-}\StringTok{ }\KeywordTok{data.frame}\NormalTok{(jml.penduduk, }\DataTypeTok{nama_kelurahan =}\NormalTok{ kelurahan)}
\end{Highlighting}
\end{Shaded}

\begin{Shaded}
\begin{Highlighting}[]
\CommentTok{# Menampilkan isi set data cek.kelurahan}
\KeywordTok{head}\NormalTok{(cek.kelurahan)}
\end{Highlighting}
\end{Shaded}

\begin{verbatim}
##      kelurahan jumlah nama_kelurahan
## 1        ANCOL  28870          Ancol
## 2        ANGKE  34663          Angke
## 3 BALE KAMBANG  32083   Bale Kambang
## 4  BALI MESTER  11290     Balimester
## 5   BAMBU APUS  28952     Bambu Apus
## 6       BANGKA  25245         Bangka
\end{verbatim}

\begin{Shaded}
\begin{Highlighting}[]
\CommentTok{# Mengecek adakah perbedaan penamaan kelurahan dengan membuat}
\CommentTok{# kolom baru bernama cek_nama yang berisi data logikal menggunakan}
\CommentTok{# menggunakan fungsi mutate() dari paket dplyr pada set data cek.kelurahan}
\NormalTok{cek.kelurahan <-}\StringTok{  }
\StringTok{  }\NormalTok{cek.kelurahan }\OperatorTok{%>%}
\StringTok{  }\KeywordTok{mutate}\NormalTok{(}\DataTypeTok{cek_nama =} \KeywordTok{ifelse}\NormalTok{(}\KeywordTok{tolower}\NormalTok{(kelurahan) }\OperatorTok{%in%}\StringTok{ }\KeywordTok{tolower}\NormalTok{(nama_kelurahan), }\StringTok{"TRUE"}\NormalTok{, }\StringTok{"FALSE"}\NormalTok{))}
\end{Highlighting}
\end{Shaded}

Operator \texttt{==} bersifat \emph{case sensitive}, artinya perbedaan
huruf besar/kecil juga menjadi faktor pengetesan logika.

Contohnya: \texttt{ifelse("ANCOL"\ ==\ "Ancol",\ "TRUE",\ "FALSE")} akan
menghasilkan output \texttt{FALSE}.

Fungsi \texttt{tolower()} yang diaplikasikan di variabel kelurahan dan
nama\_kelurahan membuat isi dari kedua variabel tersebut menjadi huruf
kecil semua sehingga pengetesan logika menjadi lebih optimal.

\begin{Shaded}
\begin{Highlighting}[]
\CommentTok{# Mengecek nama kelurahan yang berbeda dengan fungsi filter}
\NormalTok{cek.kelurahan }\OperatorTok{%>%}
\StringTok{  }\KeywordTok{select}\NormalTok{(kelurahan, cek_nama) }\OperatorTok{%>%}
\StringTok{  }\KeywordTok{filter}\NormalTok{(cek_nama }\OperatorTok{==}\StringTok{ "FALSE"}\NormalTok{)}
\end{Highlighting}
\end{Shaded}

\begin{verbatim}
##                kelurahan cek_nama
## 1            BALI MESTER    FALSE
## 2  HALIM PERDANA KUSUMAH    FALSE
## 3          HARAPAN MULIA    FALSE
## 4              KALIDERES    FALSE
## 5              KERENDANG    FALSE
## 6       KOTA BAMBU UTARA    FALSE
## 7             P. HARAPAN    FALSE
## 8              P. KELAPA    FALSE
## 9            P. PANGGANG    FALSE
## 10               P. PARI    FALSE
## 11             P. TIDUNG    FALSE
## 12        P. UNTUNG JAWA    FALSE
## 13            PAL MERIAM    FALSE
## 14              PALMERAH    FALSE
\end{verbatim}

Ternyata ada 14 nama kelurahan yang berbeda. Untuk memperbaiki tulisan
nama kelurahan, harus dibandingkan dulu versi nama kelurahan dari kedua
set data tersebut.

\begin{Shaded}
\begin{Highlighting}[]
\CommentTok{# Mengecek nama kelurahan yang berbeda dengan fungsi filter}
\NormalTok{cek.kelurahan }\OperatorTok{%>%}
\StringTok{  }\KeywordTok{select}\NormalTok{(nama_kelurahan) }\OperatorTok{%>%}
\StringTok{  }\KeywordTok{filter}\NormalTok{(}\KeywordTok{grepl}\NormalTok{(}\StringTok{"mester|halim|harapan|deres|endang|bambu|}
\StringTok{               |kelapa|panggang|pari|tidung|untung|meriam|merah"}\NormalTok{,    }
\NormalTok{               nama_kelurahan, }\DataTypeTok{ignore.case =} \OtherTok{TRUE}\NormalTok{))}
\end{Highlighting}
\end{Shaded}

\begin{verbatim}
##          nama_kelurahan
## 1            Balimester
## 2            Bambu Apus
## 3  Halim Perdana Kusuma
## 4         Harapan Mulya
## 5            Kali Deres
## 6          Kebon Kelapa
## 7            Kelapa Dua
## 8      Kelapa Dua Wetan
## 9   Kelapa Gading Barat
## 10  Kelapa Gading Timur
## 11   Kota Bambu Selatan
## 12      Kotabambu Utara
## 13             Krendang
## 14            Pal Merah
## 15            Palmeriam
## 16         Pondok Bambu
## 17        Pondok Kelapa
## 18        PULAU HARAPAN
## 19         PULAU KELAPA
## 20       PULAU PANGGANG
## 21           PULAU PARI
## 22         PULAU TIDUNG
## 23    PULAU UNTUNG JAWA
## 24         Sungai Bambu
\end{verbatim}

Dari hasil di atas, terdapat 14 perbedaan penamaan kelurahan di antara
dua set data sebagai berikut:

\begin{enumerate}
\def\labelenumi{\arabic{enumi}.}
\item
  bali mester - balimester
\item
  halim perdana kusumah - halim perdana kusuma
\item
  harapan mulia - harapan mulya
\item
  kalideres - kali deres
\item
  kerendang - krendang
\item
  kota bambu utara - kotabambu utara
\item
  \begin{enumerate}
  \def\labelenumii{\alph{enumii}.}
  \setcounter{enumii}{15}
  \tightlist
  \item
    harapan - pulau harapan
  \end{enumerate}
\item
  \begin{enumerate}
  \def\labelenumii{\alph{enumii}.}
  \setcounter{enumii}{15}
  \tightlist
  \item
    kelapa - pulau kelapa
  \end{enumerate}
\item
  \begin{enumerate}
  \def\labelenumii{\alph{enumii}.}
  \setcounter{enumii}{15}
  \tightlist
  \item
    panggang - pulau panggang
  \end{enumerate}
\item
  \begin{enumerate}
  \def\labelenumii{\alph{enumii}.}
  \setcounter{enumii}{15}
  \tightlist
  \item
    pari - pulau pari
  \end{enumerate}
\item
  \begin{enumerate}
  \def\labelenumii{\alph{enumii}.}
  \setcounter{enumii}{15}
  \tightlist
  \item
    tidung - pulau tidung
  \end{enumerate}
\item
  \begin{enumerate}
  \def\labelenumii{\alph{enumii}.}
  \setcounter{enumii}{15}
  \tightlist
  \item
    untung jawa - pulau untung jawa
  \end{enumerate}
\item
  pal meriam - palmeriam
\item
  palmerah - pal merah
\end{enumerate}

Selanjutnya, mengubah nama kelurahan di salah satu set data. Untuk
contoh kali ini, nama kelurahan pada set data \texttt{jml.penduduk} akan
diubah mengikuti penamaan kelurahan set data \texttt{kelurahan.jkt}.

Setelah penamaan kelurahan sudah sama persis, penggabungan kedua set
data ini bisa dilakukan.

\begin{Shaded}
\begin{Highlighting}[]
\CommentTok{# Mengubah nama kelurahan menjadi huruf kecil semua pada set data}
\CommentTok{# jml.penduduk dan kelurahan.jkt untuk memudahkan manipulasi string}
\NormalTok{jml.penduduk}\OperatorTok{$}\NormalTok{kelurahan <-}\StringTok{ }\KeywordTok{as.factor}\NormalTok{(}\KeywordTok{tolower}\NormalTok{(jml.penduduk}\OperatorTok{$}\NormalTok{kelurahan))}
\NormalTok{kelurahan.jkt}\OperatorTok{$}\NormalTok{nama_kelurahan <-}\StringTok{ }\KeywordTok{as.factor}\NormalTok{(}\KeywordTok{tolower}\NormalTok{(kelurahan.jkt}\OperatorTok{$}\NormalTok{nama_kelurahan))}

\CommentTok{# Mengubah penamaan kelurahan pada set data jml.penduduk}
\CommentTok{# mengikuti peamaan kelurahan set data kelurahan.jkt}
\CommentTok{# menggunakan fungsi fct_recode}
\NormalTok{jml.penduduk <-}\StringTok{ }
\StringTok{  }\KeywordTok{mutate}\NormalTok{(jml.penduduk, }\DataTypeTok{kelurahan =} 
           \KeywordTok{fct_recode}\NormalTok{(kelurahan,}
             \StringTok{"balimester"}\NormalTok{ =}\StringTok{ "bali mester"}\NormalTok{,}
             \StringTok{"halim perdana kusuma"}\NormalTok{ =}\StringTok{ "halim perdana kusumah"}\NormalTok{,}
             \StringTok{"harapan mulya"}\NormalTok{ =}\StringTok{ "harapan mulia"}\NormalTok{,}
             \StringTok{"kali deres"}\NormalTok{ =}\StringTok{ "kalideres"}\NormalTok{,}
             \StringTok{"krendang"}\NormalTok{ =}\StringTok{ "kerendang"}\NormalTok{,}
             \StringTok{"kotabambu utara"}\NormalTok{ =}\StringTok{ "kota bambu utara"}\NormalTok{,}
             \StringTok{"pulau harapan"}\NormalTok{ =}\StringTok{ "p. harapan"}\NormalTok{,}
             \StringTok{"pulau kelapa"}\NormalTok{ =}\StringTok{ "p. kelapa"}\NormalTok{,}
             \StringTok{"pulau panggang"}\NormalTok{ =}\StringTok{ "p. panggang"}\NormalTok{,}
             \StringTok{"pulau pari"}\NormalTok{ =}\StringTok{ "p. pari"}\NormalTok{,}
             \StringTok{"pulau tidung"}\NormalTok{ =}\StringTok{ "p. tidung"}\NormalTok{,}
             \StringTok{"pulau untung jawa"}\NormalTok{ =}\StringTok{ "p. untung jawa"}\NormalTok{,}
             \StringTok{"palmeriam"}\NormalTok{ =}\StringTok{ "pal meriam"}\NormalTok{,}
             \StringTok{"pal merah"}\NormalTok{ =}\StringTok{ "palmerah"} 
\NormalTok{  ))}
\end{Highlighting}
\end{Shaded}

String di sebelah kiri \texttt{=} adalah nama baru, sedangkan di sebelah
kanan \texttt{=} adalah nama lama.

Fungsi \texttt{fct\_recode()} pada dasarnya sama dengan fungsi
\texttt{gsub()}. Perbedaannya, mengubah string dengan menggunakan
\texttt{fct\_recode()} dilakukan pada variabel faktor, sedangkan
\texttt{gsub()} pada tipe variabel karakter.

Jika \texttt{gsub()} digunakan untuk memanipulasi string pada tipe
variabel faktor, maka ketika diaplikasikan variabel tersebut akan
berubah menjadi tipe variabel karakter.

\begin{Shaded}
\begin{Highlighting}[]
\CommentTok{# Menggabungkan set data jml.penduduk ke set data kelurahan.jkt}
\CommentTok{# menggunakan fungsi left_join dari paket dplyr}
\NormalTok{kelurahan.jkt <-}
\StringTok{  }\KeywordTok{left_join}\NormalTok{(kelurahan.jkt, jml.penduduk, }\DataTypeTok{by =} \KeywordTok{c}\NormalTok{(}\StringTok{"nama_kelurahan"}\NormalTok{ =}\StringTok{ "kelurahan"}\NormalTok{))}
\end{Highlighting}
\end{Shaded}

\begin{Shaded}
\begin{Highlighting}[]
\CommentTok{# Memvisualisasikan peta jumlah penduduk Jakarta}
\NormalTok{kelurahan.jkt }\OperatorTok{%>%}
\StringTok{  }\KeywordTok{filter}\NormalTok{(}\OperatorTok{!}\NormalTok{nama_kecamatan }\OperatorTok{%in%}\StringTok{ }\KeywordTok{c}\NormalTok{(}\StringTok{"KEPULAUAN SERIBU UTARA"}\NormalTok{, }\StringTok{"KEPULAUAN SERIBU SELATAN"}\NormalTok{)) }\OperatorTok{%>%}
\StringTok{  }\KeywordTok{ggplot}\NormalTok{() }\OperatorTok{+}\StringTok{ }
\StringTok{  }\KeywordTok{geom_sf}\NormalTok{(}\KeywordTok{aes}\NormalTok{(}\DataTypeTok{fill =}\NormalTok{ jumlah), }\DataTypeTok{color =} \StringTok{"gray75"}\NormalTok{) }\OperatorTok{+}
\StringTok{  }\KeywordTok{scale_fill_viridis_c}\NormalTok{() }\OperatorTok{+}
\StringTok{  }\KeywordTok{ggtitle}\NormalTok{(}\StringTok{"Jumlah Penduduk DKI Jakarta 2017"}\NormalTok{, }
          \DataTypeTok{subtitle =} \StringTok{"sumber set data: data.jakarta.go.id; referensi kode: datascience.or.id"}\NormalTok{) }\OperatorTok{+}
\StringTok{  }\KeywordTok{theme}\NormalTok{(}\DataTypeTok{panel.background =} \KeywordTok{element_blank}\NormalTok{()) }\OperatorTok{+}
\StringTok{  }\KeywordTok{theme}\NormalTok{(}\DataTypeTok{axis.title =} \KeywordTok{element_blank}\NormalTok{()) }\OperatorTok{+}
\StringTok{  }\KeywordTok{theme}\NormalTok{(}\DataTypeTok{axis.text =} \KeywordTok{element_blank}\NormalTok{()) }\OperatorTok{+}
\StringTok{  }\KeywordTok{theme}\NormalTok{(}\DataTypeTok{axis.ticks =} \KeywordTok{element_blank}\NormalTok{()) }\OperatorTok{+}
\StringTok{  }\KeywordTok{theme}\NormalTok{(}\DataTypeTok{legend.title =} \KeywordTok{element_blank}\NormalTok{())}
\end{Highlighting}
\end{Shaded}

\includegraphics{ebook_files/figure-latex/unnamed-chunk-104-1.pdf}

Jika Anda ingin menampilkan salah satu kota di DKI Jakarta, tinggal
menambahkan saringan set data untuk kota terkait. Contoh berikut
menampilkan kota Jakarta Barat.

\begin{Shaded}
\begin{Highlighting}[]
\CommentTok{# Memvisualisasikan peta jumlah penduduk Jakarta Barat}
\NormalTok{kelurahan.jkt }\OperatorTok{%>%}
\StringTok{  }\KeywordTok{filter}\NormalTok{(}\OperatorTok{!}\NormalTok{nama_kecamatan }\OperatorTok{%in%}\StringTok{ }\KeywordTok{c}\NormalTok{(}\StringTok{"KEPULAUAN SERIBU UTARA"}\NormalTok{, }\StringTok{"KEPULAUAN SERIBU SELATAN"}\NormalTok{) }\OperatorTok{&}\StringTok{ }
\StringTok{          }\NormalTok{nama_kota }\OperatorTok{==}\StringTok{ "Jakarta Barat"}\NormalTok{) }\OperatorTok{%>%}
\StringTok{  }\KeywordTok{ggplot}\NormalTok{() }\OperatorTok{+}\StringTok{ }
\StringTok{  }\KeywordTok{geom_sf}\NormalTok{(}\KeywordTok{aes}\NormalTok{(}\DataTypeTok{fill =}\NormalTok{ jumlah), }\DataTypeTok{color =} \StringTok{"gray75"}\NormalTok{) }\OperatorTok{+}
\StringTok{  }\KeywordTok{scale_fill_viridis_c}\NormalTok{() }\OperatorTok{+}
\StringTok{  }\KeywordTok{ggtitle}\NormalTok{(}\StringTok{"Jumlah Penduduk Jakarta Barat 2017"}\NormalTok{,}
          \DataTypeTok{subtitle =} \StringTok{"sumber set data: data.jakarta.go.id; referensi kode: datascience.or.id"}\NormalTok{) }\OperatorTok{+}
\StringTok{  }\KeywordTok{theme}\NormalTok{(}\DataTypeTok{panel.background =} \KeywordTok{element_blank}\NormalTok{()) }\OperatorTok{+}
\StringTok{  }\KeywordTok{theme}\NormalTok{(}\DataTypeTok{axis.title =} \KeywordTok{element_blank}\NormalTok{()) }\OperatorTok{+}
\StringTok{  }\KeywordTok{theme}\NormalTok{(}\DataTypeTok{axis.text =} \KeywordTok{element_blank}\NormalTok{()) }\OperatorTok{+}
\StringTok{  }\KeywordTok{theme}\NormalTok{(}\DataTypeTok{axis.ticks =} \KeywordTok{element_blank}\NormalTok{()) }\OperatorTok{+}
\StringTok{  }\KeywordTok{theme}\NormalTok{(}\DataTypeTok{legend.title =} \KeywordTok{element_blank}\NormalTok{())}
\end{Highlighting}
\end{Shaded}

\includegraphics{ebook_files/figure-latex/unnamed-chunk-105-1.pdf}

\section{\texorpdfstring{\emph{Insight}}{Insight}}\label{insight}

Beberapa \emph{insight} atau informasi yang bisa didapat dari hasil
eksplorasi di atas antara lain sebagai berikut:

\begin{enumerate}
\def\labelenumi{\arabic{enumi}.}
\tightlist
\item
  Jumlah penduduk terbanyak ada di kecamatan Cengkareng, sedangkan untuk
  kategori kota ada di Jakarta Timur.
\item
  Jumlah penduduk paling sedikit terdapat di Kepulauan Seribu.
\item
  DKI Jakarta dalam angka jumlah penduduk tahun 2017:

  \begin{itemize}
  \tightlist
  \item
    Total = 10.348.570.
  \item
    Kecamatan terbanyak, Cengkareng = 513.064.
  \item
    Kecamatan/kota tersedikit, Kepulauan Seribu = 27.591.
  \item
    Kota terbanyak, Jakarta Timur = 2.946.926.
  \end{itemize}
\end{enumerate}

\section{Referensi}\label{referensi}

Berikut referensi yang bisa Anda jadikan panduan untuk lebih cepat
menguasai R:

\begin{itemize}
\item
  \href{https://www.rstudio.com/resources/cheatsheets/}{\emph{Cheat
  Sheet} RStudio}. \emph{Cheat Sheet} atau contekan yang dirangkum dalam
  2-3 lembar ini membantu dan mempermudah Anda untuk menavigasi
  fungsi-fungsi yang diperlukan ketika proses eksplorasi data.
\item
  \href{https://www.datacamp.com/courses/free-introduction-to-r}{Kursus
  online gratis dari DataCamp}. Tidak hanya belajar teori dan perintah
  dasar pemrograman R, Anda juga akan disuguhi dan dilatih dengan banyak
  praktek penulisan kode. Setiap bagian akan ada instruksi juga kuis
  yang harus Anda jawab dengan menuliskan kode yang sesuai.
\item
  Googling. Anda juga bisa mendapatkan referensi melalui Google dengan
  mengetikkan kata kunci secara spesifik terkait permasalahan tentang R
  yang Anda hadapi.
\end{itemize}

Akhir kata, semoga buku yang ditulis secara amatir ini dapat membawa
banyak manfaat bagi Anda.

Terima kasih sudah membaca buku ini sampai akhir.

dan..

\textbf{SELAMAT BELAJAR}

\begin{quote}
``Anyone who stops learning is old, whether at twenty or eighty. Anyone
who keeps learning stays young.''
\end{quote}

\begin{quote}
―- \textbf{Henry Ford}
\end{quote}

\bibliography{book.bib,packages.bib}


\end{document}
